\NormalFont\ShortTitle{Matei}

{\MTB Vestea cea bună conform

\par }{\MT Matei


\par }\ChapOne{1}{\PP \VerseOne{1}Cartea genealogiei lui Isus Hristos,\FTNT{*}{{\FR 1:1 }{\FT{Mesia (ebraică) și Hristos (greacă) înseamnă amândouă “Unsul”.}}} fiul lui David, fiul lui Avraam.


\par }{\PP \VS{2}Avraam a devenit tatăl lui Isaac. Isaac a devenit tatăl lui Iacov. Iacov a devenit tatăl lui Iuda și al fraților săi.

\VS{3}Iuda a devenit tatăl lui Pereț și al lui Zerah prin Tamar. Perez a devenit tatăl lui Hezron. Hezron a devenit tatăl lui Berbec.

\VS{4}Berbec a devenit tatăl lui Aminadab. Amminadab a devenit tatăl lui Nahșon. Nahșon a devenit tatăl lui Salmon.

\VS{5}Salmon a devenit tatăl lui Boaz prin Rahab. Boaz a devenit tatăl lui Obed prin Rut. Obed a devenit tatăl lui Isai.

\VS{6}Isai a devenit tatăl regelui David. Regele\FTNT{†}{{\FR 1:6 }{\FT{NU omite “regele”.}}} David a devenit tatăl lui Solomon prin cea care fusese soția lui Urie.

\VS{7}Solomon a devenit tatăl lui Roboam. Roboam a devenit tatăl lui Abia. Abiia a devenit tatăl lui Asa.

\VS{8}Asa a devenit tatăl lui Iosafat. Iosafat a devenit tatăl lui Ioram. Ioram a devenit tatăl lui Ozia.

\VS{9}Ozia a devenit tatăl lui Iotam. Iotam a devenit tatăl lui Ahaz. Ahaz a devenit tatăl lui Ezechia.

\VS{10}Ezechia a devenit tatăl lui Manase. Manase a devenit tatăl lui Amon. Amon a devenit tatăl lui Iosia.

\VS{11}Iosia a devenit tatăl lui Ieconia și al fraților săi în timpul exilului în Babilon.


\par }{\PP \VS{12}După exilul în Babilon, Ieconia a devenit tatăl lui Șealtiel. Șealtiel a devenit tatăl lui Zorobabel.

\VS{13}Zorobabel a devenit tatăl lui Abiud. Abiud a devenit tatăl lui Eliachim. Eliachim a devenit tatăl lui Azor.

\VS{14}Azor a devenit tatăl lui Țadoc. Țadoc a devenit tatăl lui Achim. Achim a devenit tatăl lui Eliud.

\VS{15}Eliud a devenit tatăl lui Eleazar. Eleazar a devenit tatăl lui Matan. Mathan a devenit tatăl lui Iacov.

\VS{16}Iacov a devenit tatăl lui Iosif, soțul Mariei, din care s-a născut Isus,
\FTNT{‡}{{\FR 1:16 }{\FT{“Isus” înseamnă “mântuire”.}}}care se numește Hristos.


\par }{\PP \VS{17}Astfel, toate generațiile de la Avraam până la David sunt de paisprezece generații; de la David până la exilul în Babilon, paisprezece generații; și de la exilul în Babilon până la Hristos, paisprezece generații.


\par }{\PP \VS{18}Iată cum a fost nașterea lui Isus Hristos: După ce mama Lui, Maria, a fost logodită cu Iosif, înainte de a fi împreună, a fost găsită însărcinată de Duhul Sfânt.

\VS{19}Iosif, soțul ei, care era un om drept și care nu voia să facă din ea un exemplu public, a vrut să o renege în secret.

\VS{20}Dar când se gândea la aceste lucruri, iată că
\FTNT{§}{{\FR 1:20 }{\FT{“Iată”, de la “ἰδοὺ”, înseamnă a privi, a lua aminte, a observa, a vedea sau a privi. Este adesea folosit ca o interjecție.}}}un înger al Domnului i s-a arătat în vis și i-a zis: “Iosif, fiul lui David, nu te teme să iei de soție pe Maria, căci ceea ce este conceput în ea este de la Duhul Sfânt.

\VS{21}Ea va naște un fiu. Să-i pui numele Isus,
\FTNT{*}{{\FR 1:21 }{\FT{“Isus” înseamnă “mântuire”.}}}căci El este cel care va mântui poporul Său de păcatele sale.”


\par }{\PP \VS{22}Și toate acestea s-au întâmplat ca să se împlinească ce a spus Domnul prin proorocul care a zis,


\par }{\Q \VS{23}“Iată, fecioara va rămâne însărcinată,

\par }{\QB și va da naștere unui fiu.

\par }{\Q Îi vor pune numele Emanuel”.

\par }{\QB care este, fiind interpretat, “Dumnezeu cu noi”.
\FTNT{†}{{\FR 1:23 }{\FT{Isaia 7:14}}}


\par }{\PP \VS{24}Iosif s-a sculat din somn, a făcut cum îi poruncise îngerul Domnului și a luat la el pe nevastă-sa,

\VS{25}și n-a cunoscut-o până ce n-a născut pe fiul ei cel întâi născut. I-a pus numele Isus.



\par }\Chap{2}{\PP \VerseOne{1}Pe când se născuse Isus în Betleemul Iudeii, în zilele împăratului Irod, iată că au venit la Ierusalim niște înțelepți de
\FTNT{*}{{\FR 2:1 }{\FT{Cuvântul pentru “înțelepți” (magoi) poate însemna și învățători, oameni de știință, medici, astrologi, clarvăzători, interpreți de vise sau vrăjitori.}}}la răsărit, care ziceau:

\VS{2}“Unde este Cel ce S-a născut Împărat al Iudeilor? Căci am văzut steaua Lui la răsărit și am venit să ne închinăm Lui”.

\VS{3}Regele Irod, când a auzit aceasta, s-a tulburat și tot Ierusalimul împreună cu el.

\VS{4}Adunând pe toți preoții cei mai de seamă și pe cărturarii poporului, i-a întrebat unde se va naște Hristosul.

\VS{5}Ei i-au răspuns: “În Betleemul Iudeii, căci așa este scris prin profet,


\par }{\Q \VS{6}“Tu, Betleem, țară a lui Iuda,

\par }{\QB nu sunt cu nimic mai prejos printre căpeteniile lui Iuda;

\par }{\Q pentru că din tine va ieși un guvernator

\par }{\QB care va păstori poporul Meu, Israel.””
\FTNT{†}{{\FR 2:6 }{\FT{Mica 5:2}}}


\par }{\PP \VS{7}Irod a chemat pe ascuns pe înțelepți și a aflat de la ei la ce oră a apărut steaua.

\VS{8}I-a trimis la Betleem și le-a zis: “Mergeți și căutați cu grijă pruncul. Când îl veți găsi, aduceți-mi vestea, ca să vin și eu să mă închin lui.”


\par }{\PP \VS{9}Ei, după ce au auzit pe rege, au plecat; și iată că steaua pe care o văzuseră la răsărit mergea înaintea lor, până ce a venit și s-a oprit deasupra locului unde era pruncul.

\VS{10}Când au văzut steaua, s-au bucurat cu bucurie nespus de mare.

\VS{11}Au intrat în casă și au văzut Pruncul cu Maria, mama Lui, și au căzut jos și I s-au închinat. Deschizându-și comorile, i-au oferit daruri: aur, tămâie și smirnă.

\VS{12}Fiind avertizați în vis să nu se întoarcă la Irod, s-au întors în țara lor pe un alt drum.


\par }{\PP \VS{13}După ce au plecat, iată că un înger al Domnului s-a arătat în vis lui Iosif și i-a zis: “Scoală-te, ia Pruncul și pe mama lui și fugi în Egipt și rămâi acolo până ce-ți voi spune, căci Irod va căuta Pruncul ca să-l nimicească.”


\par }{\PP \VS{14}S-a sculat, a luat noaptea pe prunc și pe mama lui și a plecat în Egipt,

\VS{15}și a stat acolo până la moartea lui Irod, ca să se împlinească ce spusese Domnul prin proorocul care zicea: “Din Egipt am chemat pe fiul Meu”.
\FTNT{‡}{{\FR 2:15 }{\FT{Osea 11:1}}}


\par }{\PP \VS{16}Atunci Irod, văzând că este batjocorit de magi, s-a mâniat foarte tare și a trimis să ucidă pe toți copiii de parte bărbătească din Betleem și din toată ținutul din jur, de la doi ani în jos, după timpul exact pe care-l aflase de la magi.

\VS{17}Atunci s-a împlinit ceea ce fusese spus prin profetul Ieremia, care zicea așa


\par }{\Q \VS{18}“S-a auzit o voce în Rama,

\par }{\QB jale, plâns și jale mare,

\par }{\Q Rahela plângând pentru copiii ei;

\par }{\QB ea nu ar fi fost consolată,

\par }{\QB pentru că ei nu mai sunt.”
\FTNT{§}{{\FR 2:18 }{\FT{Ieremia 31:15}}}


\par }{\PP \VS{19}După ce a murit Irod, iată că un înger al Domnului s-a arătat în vis lui Iosif în Egipt, zicând:

\VS{20}“Scoală-te, ia Pruncul și pe mama lui și du-te în țara lui Israel, căci cei care căutau viața Pruncului au murit.”


\par }{\PP \VS{21}Și s-a sculat, a luat pruncul și pe mama lui și a venit în țara lui Israel.

\VS{22}Dar, când a auzit că Arhelau domnea în Iudeea în locul tatălui său, Irod, s-a temut să meargă acolo. Fiind avertizat în vis, s-a retras în regiunea Galileii

\VS{23}și a venit și a locuit într-o cetate numită Nazaret, ca să se împlinească ceea ce s-a spus prin profeți, că el va fi numit Nazarinean.



\par }\Chap{3}{\PP \VerseOne{1}În zilele acelea, a venit Ioan Botezătorul și propovăduia în pustiul Iudeii, zicând:

\VS{2}“Pocăiți-vă, căci Împărăția cerurilor este aproape!”

\VS{3}Căci acesta este cel despre care s-a vorbit prin profetul Isaia, care a spus

\par }{\Q “Glasul cuiva care strigă în pustiu,

\par }{\QB pregătiți calea Domnului!

\par }{\QB Fă-i cărările drepte!”
\FTNT{*}{{\FR 3:3 }{\FT{Isaia 40:3}}}


\par }{\PP \VS{4}Ioan purta haine de păr de cămilă și o curea de piele la brâu. Hrana lui era lăcuste și miere sălbatică.

\VS{5}Atunci a ieșit la el popor din Ierusalim, din toată Iudeea și din toată regiunea din jurul Iordanului.

\VS{6}Au fost botezați de el în Iordan, mărturisindu-și păcatele.


\par }{\PP \VS{7}Dar când a văzut pe mulți dintre farisei și saduchei venind la botezul Său, le-a zis: “Urmași ai viperelor, cine v-a avertizat să fugiți de mânia viitoare?

\VS{8}De aceea, produceți roade demne de pocăință!

\VS{9}Nu vă gândiți în sinea voastră: “Îl avem ca tată pe Avraam!”, căci vă spun că Dumnezeu poate să-i ridice lui Avraam copii din aceste pietre.

\VS{10}Chiar acum securea stă la rădăcina copacilor. De aceea, orice copac care nu produce fructe bune este tăiat și aruncat în foc.


\par }{\PP \VS{11}Eu vă botez în apă pentru pocăință, dar Cel ce vine după mine este mai puternic decât mine, și nu sunt vrednic să port sandalele lui. El vă va boteza în Duhul Sfânt.
\FTNT{†}{{\FR 3:11 }{\FT{TR și NU adaugă “și cu foc”
}}}

\VS{12}Furca lui de vânat este în mâna lui și își va curăța temeinic aria de treierat. Își va strânge grâul în hambar, dar neghina o va arde cu foc nestins.”


\par }{\PP \VS{13}Atunci Isus a venit din Galileea la Iordan, la Ioan, ca să fie botezat de el.

\VS{14}Dar Ioan ar fi vrut să-L împiedice, zicând: “Eu am nevoie să fiu botezat de tine, și tu vii la mine?”


\par }{\PP \VS{15}Isus, răspunzând, i-a zis: “{\WJ{Îngăduie acum, căci așa se cuvine să împlinim toată dreptatea.” Atunci el i-a permis.
}}


\par }{\PP \VS{16}Isus, după ce a fost botezat, s-a suit din apă și iată că i s-au deschis cerurile. A văzut Duhul lui Dumnezeu coborând ca un porumbel și venind peste el.

\VS{17}Și iată, un glas din ceruri a zis: “Acesta este Fiul Meu preaiubit, cu care sunt foarte mulțumit.”



\par }\Chap{4}{\PP \VerseOne{1}Isus a fost dus de Duhul Sfânt în pustiu, ca să fie ispitit de diavol.

\VS{2}După ce a postit patruzeci de zile și patruzeci de nopți, I-a fost foame după aceea.

\VS{3}Ispititorul a venit și i-a zis: “Dacă ești Fiul lui Dumnezeu, poruncește ca aceste pietre să se transforme în pâine”.


\par }{\PP \VS{4}Iar el a răspuns: “{\WJ{Este scris: “Nu numai cu pâine va trăi omul, ci cu orice cuvânt care iese din gura lui Dumnezeu.””
}}


\par }{\PP \VS{5}Atunci diavolul l-a dus în cetatea sfântă. L-a așezat pe vârful templului

\VS{6}și i-a zis: “Dacă ești Fiul lui Dumnezeu, aruncă-te jos, căci este scris,

\par }{\Q “El va porunci îngerilor Săi cu privire la voi” și,

\par }{\Q 'Pe mâinile lor te vor purta,

\par }{\QB ca să nu te lovești cu piciorul de o piatră.””


\par }{\PP \VS{7}Isus i-a zis: “{\WJ{Iarăși este scris: “Să nu pui la încercare pe Domnul, Dumnezeul tău”.”
}}


\par }{\PP \VS{8}Diavolul l-a dus iarăși pe un munte foarte înalt și i-a arătat toate împărățiile lumii și slava lor.

\VS{9}El i-a spus: “Îți voi da toate aceste lucruri, dacă vei cădea jos și te vei închina mie.”


\par }{\PP \VS{10}Atunci Isus i-a zis: “{\WJ{Du-te de aici, Satano! Căci este scris: “Să te închini Domnului, Dumnezeului tău, și numai lui să-i slujești”.”
}}


\par }{\PP \VS{11}Atunci diavolul L-a lăsat și iată că au venit îngeri și I-au slujit.


\par }{\PP \VS{12}Când a auzit Isus că Ioan fusese predat, s-a retras în Galileea.

\VS{13}A plecat din Nazaret, a venit și a locuit în Capernaum, care este lângă mare, în ținutul lui Zabulon și Neftali,

\VS{14}ca să se împlinească ceea ce fusese spus prin profetul Isaia, care zicea:,


\par }{\Q \VS{15}“Țara lui Zabulon și țara lui Neftali,

\par }{\QB spre mare, dincolo de Iordan,

\par }{\QB Galileea neamurilor,


\par }{\Q \VS{16}poporul care ședea în întuneric a văzut o lumină mare;

\par }{\QB celor care stăteau în regiunea și în umbra morții,

\par }{\QB pentru ei a răsărit lumina.”


\par }{\PP \VS{17}Din vremea aceea, Isus a început să propovăduiască și să zică:
{\WJ{“Pocăiți-vă! Căci Împărăția Cerurilor este aproape”.
}}


\par }{\PP \VS{18}Și mergând pe malul mării Galileii, a văzut doi frați: Simon, numit Petru, și Andrei, fratele său, aruncând o plasă în mare, căci erau pescari.

\VS{19}El le-a zis:
{\WJ{“Veniți după Mine și vă voi face pescari de oameni.”
}}


\par }{\PP \VS{20}Și îndată și-au lăsat plasele și au mers după El.

\VS{21}Mergând mai departe de acolo, a văzut alți doi frați, Iacov, fiul lui Zebedei, și Ioan, fratele său, în barcă, împreună cu Zebedei, tatăl lor, reparându-și plasele. El i-a chemat.

\VS{22}Ei au lăsat îndată barca și pe tatăl lor și l-au urmat.


\par }{\PP \VS{23}Isus străbătea toată Galileea, învățând în sinagogile lor, propovăduind vestea cea bună a Împărăției și vindecând orice boală și orice neputință în popor.

\VS{24}Vestea despre El s-a răspândit în toată Siria. Îi aduceau la el pe toți cei care erau bolnavi, atinși de diferite boli și chinuri, posedați de demoni, epileptici și paralitici; și el îi vindeca.

\VS{25}L-au urmat mari mulțimi din Galileea, din Decapole, din Ierusalim, din Iudeea și de dincolo de Iordan.



\par }\Chap{5}{\PP \VerseOne{1}Văzând mulțimile, S-a suit pe munte. După ce a șezut, discipolii Lui au venit la El.

\VS{2}El a deschis gura și i-a învățat, zicând:,


\par }{\Q \VS{3}{\WJ{“Ferice de cei săraci cu duhul!
}}

\par }{\QB {\WJ{pentru că a lor este Împărăția Cerurilor.
}}


\par }{\Q \VS{4}{\WJ{Ferice de cei ce plâng!
}}

\par }{\QB {\WJ{pentru că vor fi mângâiați.
}}


\par }{\Q \VS{5}{\WJ{Ferice de cei blânzi!
}}

\par }{\QB {\WJ{căci ei vor moșteni pământul.
}}


\par }{\Q \VS{6}{\WJ{Ferice de cei ce flămânzesc și însetează după dreptate}},

\par }{\QB {\WJ{pentru că vor fi umplute.
}}


\par }{\Q \VS{7}{\WJ{Ferice de cei milostivi!
}}

\par }{\QB {\WJ{pentru că ei vor obține îndurare.
}}


\par }{\Q \VS{8}{\WJ{Ferice de cei cu inima curată!
}}

\par }{\QB {\WJ{pentru că ei îl vor vedea pe Dumnezeu.
}}


\par }{\Q \VS{9}{\WJ{Fericiți făcătorii de pace!
}}

\par }{\QB {\WJ{pentru că ei se vor numi copii ai lui Dumnezeu.
}}


\par }{\Q \VS{10}{\WJ{Ferice de cei ce au fost prigoniți pentru dreptate}},

\par }{\QB {\WJ{pentru că a lor este Împărăția Cerurilor.
}}


\par }{\PP \VS{11}{\WJ{Ferice de voi, când vă vor ocărî, vă vor prigoni și vor zice tot felul de lucruri rele împotriva voastră, pe nedrept, din pricina Mea.
}}

\VS{12}{\WJ{Bucurați-vă și vă veseliți, căci mare este răsplata voastră în ceruri. Căci așa au prigonit ei pe profeții care au fost înaintea voastră.
}}


\par }{\PP \VS{13}{\WJ{Voi sunteți sarea pământului; dar dacă sarea și-a pierdut gustul, cu ce se va săra? Atunci nu mai este bună la nimic, decât să fie aruncată afară și călcată în picioare de oameni.
}}


\par }{\PP \VS{14}{\WJ{Tu ești lumina lumii. O cetate situată pe un deal nu poate fi ascunsă.
}}

\VS{15}{\WJ{Nici nu aprindeți o lampă și nu o puneți sub un coș de măsurat, ci pe un stativ, și ea luminează tuturor celor care sunt în casă.
}}

\VS{16}{\WJ{Tot așa, lumina voastră să strălucească înaintea oamenilor, ca să vadă faptele voastre bune și să slăvească pe Tatăl vostru care este în ceruri.
}}


\par }{\PP \VS{17}{\WJ{“Să nu credeți că am venit să distrug Legea și proorocii. Nu am venit să distrug, ci să împlinesc.
}}

\VS{18}{\WJ{Căci, cu siguranță, vă spun că, până când vor trece cerul și pământul, nici cea mai mică literă și nici cea mai mică lovitură de peniță nu va trece în vreun fel din Lege, până când toate lucrurile vor fi împlinite.
}}

\VS{19}{\WJ{De aceea, oricine va călca una dintre aceste porunci mai mici și va învăța pe alții să facă așa, va fi numit cel mai mic în Împărăția cerurilor; dar oricine le va împlini și le va învăța va fi numit mare în Împărăția cerurilor.
}}

\VS{20}{\WJ{Căci vă spun că, dacă dreptatea voastră nu o va întrece pe cea a cărturarilor și a fariseilor, în niciun caz nu veți intra în Împărăția Cerurilor.
}}


\par }{\PP \VS{21}Ați
{\WJ{auzit că s-a spus celor din vechime: “Să nu ucizi!” și că “Oricine va ucide va fi supus judecății”.
}}

\VS{22}{\WJ{Dar Eu vă spun că oricine se mânie pe fratele său fără motiv va fi în primejdie de judecată. Oricine îi va spune fratelui său: “Raca!}}” va
{\WJ{fi în pericol de consiliu. Oricine îi va spune: “Prostule!” va fi în pericol de focul gheenei.
}}


\par }{\PP \VS{23}“{\WJ{Deci, dacă îți aduci darul la altar și îți aduci aminte că fratele tău are ceva împotriva ta,
}}

\VS{24}{\WJ{lasă darul tău acolo, înaintea altarului, și du-te. Împăcați-vă mai întâi cu fratele vostru și apoi veniți și aduceți darul vostru.
}}

\VS{25}{\WJ{Înțelege-te repede cu potrivnicul tău, cât timp ești cu el pe drum, ca nu cumva procurorul să te predea judecătorului, iar judecătorul să te predea ofițerului și să fii aruncat în închisoare.
}}

\VS{26}{\WJ{Mai mult ca sigur vă spun că nu veți ieși nicidecum de acolo până nu veți plăti și ultimul bănuț.
}}


\par }{\PP \VS{27}Ați
{\WJ{auzit că s-a zis: “Să nu preacurvești!”
}}

\VS{28}{\WJ{Dar Eu vă spun că oricine se uită la o femeie ca să o poftească, a preacurvit cu ea în inima lui.
}}

\VS{29}{\WJ{Dacă ochiul tău drept te face să te poticnești, smulge-l și aruncă-l de la tine. Căci este mai de folos să piară unul dintre membrele tale decât ca tot trupul tău să fie aruncat în gheenă.
}}

\VS{30}{\WJ{Dacă mâna ta dreaptă te face să te poticnești, taie-o și arunc-o departe de tine. Căci îți este mai de folos să piară unul din membrele tale, decât să fie aruncat tot trupul tău în Gheenă.
}}


\par }{\PP \VS{31}“{\WJ{S-a mai spus: “Oricine își va renega nevasta, să-i dea un act de divorț.”}}

\VS{32}{\WJ{Dar Eu vă spun că oricine își reneagă nevasta, afară de motivul desfrânării, o face adulteră; și oricine se căsătorește cu ea după ce a fost renegată, comite adulter.
}}


\par }{\PP \VS{33}“Ați
{\WJ{auzit că s-a spus celor vechi: “Să nu faceți jurăminte false, ci să împliniți jurămintele voastre față de Domnul.”}}

\VS{34}{\WJ{Dar Eu vă spun să nu jurați deloc: nici pe cer, căci este tronul lui Dumnezeu,
}}

\VS{35}nici pe
{\WJ{pământ, căci este așternutul picioarelor Lui, nici pe Ierusalim, căci este cetatea marelui Împărat.
}}

\VS{36}{\WJ{Să nu juri nici pe capul tău, căci nu poți face nici un fir de păr alb sau negru.
}}

\VS{37}{\WJ{Ci să fie “Da” al tău “Da” și “Nu” al tău “Nu”. Tot ce este mai mult decât acestea este de la cel rău.
}}


\par }{\PP \VS{38}Ați
{\WJ{auzit că s-a spus: “Ochi pentru ochi și dinte pentru dinte”.
}}

\VS{39}{\WJ{Dar Eu vă spun: Nu vă împotriviți celui rău; ci, dacă te lovește cineva peste obrazul drept, întoarce-i și celălalt obraz.
}}

\VS{40}{\WJ{Dacă cineva te dă în judecată ca să-ți ia haina, lasă-i și haina ta.
}}

\VS{41}{\WJ{Oricine vă silește să mergeți o milă, mergeți cu el două.
}}

\VS{42}{\WJ{Dă-i celui care îți cere și nu refuza pe cel care vrea să se împrumute de la tine.
}}


\par }{\PP \VS{43}Ați
{\WJ{auzit că s-a spus: “Să iubești pe aproapele tău și să urăști pe vrăjmașul tău.
}}

\VS{44}{\WJ{Dar Eu vă spun: iubiți pe vrăjmașii voștri, binecuvântați pe cei ce vă blestemă, faceți bine celor ce vă urăsc și rugați-vă pentru cei ce vă maltratează și vă prigonesc,
}}

\VS{45}{\WJ{ca să fiți copii ai Tatălui vostru care este în ceruri. Căci El face să răsară soarele Său peste cei răi și peste cei buni și trimite ploaie peste cei drepți și peste cei nedrepți.
}}

\VS{46}{\WJ{Căci, dacă iubiți pe cei ce vă iubesc, ce răsplată aveți voi? Oare nici măcar vameșii nu fac la fel?
}}

\VS{47}{\WJ{Dacă nu faci decât să-ți saluți prietenii, ce faci tu mai mult decât ceilalți? Nu fac la fel și vameșii?
}}

\VS{48}{\WJ{De aceea, voi să fiți desăvârșiți, după cum Tatăl vostru din ceruri este desăvârșit.
}}



\par }\Chap{6}{\PP \VerseOne{1}“Luați
{\WJ{seama să nu faceți milostenia voastră înaintea oamenilor, ca să fiți văzuți de ei, căci altfel nu veți avea nici o răsplată de la Tatăl vostru care este în ceruri.
}}

\VS{2}{\WJ{De aceea, când faceți fapte de milostenie, nu sunați din trâmbiță înaintea voastră, așa cum fac ipocriții în sinagogi și pe străzi, ca să obțină slavă de la oameni. Mai mult ca sigur vă spun că ei și-au primit răsplata.
}}

\VS{3}{\WJ{Dar când faci fapte de milostenie, nu lăsa ca mâna ta stângă să știe ce face mâna ta dreaptă,
}}

\VS{4}pentru ca
{\WJ{faptele tale de milostenie să fie în taină; atunci Tatăl tău, care vede în taină, îți va răsplăti în mod deschis.
}}


\par }{\PP \VS{5}{\WJ{Când vă rugați, să nu fiți ca fățarnicii, căci lor le place să stea în picioare și să se roage în sinagogi și la colțurile străzilor, ca să fie văzuți de oameni. Mai mult ca sigur, vă spun, ei și-au primit răsplata.
}}

\VS{6}{\WJ{Dar tu, când te rogi, intră în camera ta interioară și, după ce ți-ai închis ușa, roagă-te Tatălui tău, care este în ascuns; și Tatăl tău, care vede în ascuns, îți va răsplăti în mod deschis.
}}

\VS{7}{\WJ{Când vă rugați, nu folosiți repetări deșarte, așa cum fac neamurile, căci ele cred că vor fi auzite pentru vorbăria lor multă.
}}

\VS{8}{\WJ{De aceea, nu fiți ca ei, căci Tatăl vostru știe de ce aveți nevoie, înainte ca voi să i le cereți.
}}

\VS{9}{\WJ{Rugați-vă astfel:
}}

\par }{\Q {\WJ{“Tatăl nostru care ești în ceruri, fie ca numele Tău să fie păstrat sfânt.
}}


\par }{\Q \VS{10}{\WJ{Să vină Împărăția Ta!
}}

\par }{\QB {\WJ{Facă-se voia Ta, precum în cer așa și pe pământ.
}}


\par }{\Q \VS{11}{\WJ{Dă-ne astăzi pâinea noastră cea de toate zilele.
}}


\par }{\Q \VS{12}{\WJ{Iartă-ne datoriile noastre,
}}

\par }{\QB {\WJ{așa cum și noi iertăm datornicilor noștri.
}}


\par }{\Q \VS{13}{\WJ{Nu ne duce pe noi în ispită,
}}

\par }{\QB {\WJ{ci izbăvește-ne de cel rău.
}}

\par }{\Q {\WJ{Căci a Ta este Împărăția, puterea și slava în veci. Amin.
}}


\par }{\PP \VS{14}{\WJ{Căci, dacă veți ierta oamenilor greșelile lor, și Tatăl vostru ceresc vă va ierta vouă.
}}

\VS{15}{\WJ{Dar dacă nu veți ierta oamenilor greșelile lor, nici Tatăl vostru nu vă va ierta greșelile voastre.
}}


\par }{\PP \VS{16}{\WJ{Și când postiți, să nu fiți ca fățarnicii, cu fețele triste. Căci ei își desfigurează fețele ca să fie văzuți de oameni că postesc. Mai mult ca sigur vă spun că ei și-au primit răsplata.
}}

\VS{17}{\WJ{Dar voi, când postiți, ungeți-vă capul și spălați-vă pe față,
}}

\VS{18}{\WJ{astfel încât să nu vă vadă oamenii că postiți, ci Tatăl vostru, care este în ascuns; și Tatăl vostru, care vede în ascuns, vă va răsplăti.
}}


\par }{\PP \VS{19}{\WJ{“Nu vă adunați comori pe pământ, unde molia și rugina mistuie și unde hoții pătrund și fură;
}}

\VS{20}{\WJ{ci adunați-vă comori în ceruri, unde nici molia, nici rugina nu mistuie și unde hoții nu pătrund și nu fură;
}}

\VS{21}{\WJ{căci unde este comoara voastră, acolo va fi și inima voastră.
}}


\par }{\PP \VS{22}{\WJ{“Lampa trupului este ochiul. De aceea, dacă ochiul tău este sănătos, tot corpul tău va fi plin de lumină.
}}

\VS{23}{\WJ{Dar dacă ochiul tău este rău, tot trupul tău va fi plin de întuneric. De aceea, dacă lumina care este în tine este întuneric, cât de mare este întunericul!
}}


\par }{\PP \VS{24}Nimeni
{\WJ{nu poate sluji la doi stăpâni, căci ori va urî pe unul și va iubi pe celălalt, ori va fi devotat unuia și va disprețui pe celălalt. Nu poți sluji și lui Dumnezeu și lui Mamona.
}}

\VS{25}{\WJ{De aceea vă spun: nu vă îngrijorați pentru viața voastră: ce veți mânca sau ce veți bea; și nici pentru trupul vostru, ce veți purta. Nu este viața mai mult decât hrana, iar trupul mai mult decât îmbrăcămintea?
}}

\VS{26}{\WJ{Priviți păsările cerului, care nu seamănă, nu seceră și nu adună în hambare. Tatăl vostru ceresc le hrănește. Nu sunteți voi mult mai valoroși decât ele?
}}


\par }{\PP \VS{27}{\WJ{“Cine dintre voi, dacă este îngrijorat, poate să adauge o clipă la viața lui?
}}

\VS{28}{\WJ{De ce vă îngrijorați din cauza hainelor? Gândiți-vă la crinii de pe câmp, cum cresc ei. Ei nu se trudesc, nici nu se învârt,
}}

\VS{29}și
{\WJ{totuși vă spun că nici măcar Solomon, în toată gloria lui, nu era îmbrăcat ca unul dintre aceștia.
}}

\VS{30}{\WJ{Dar dacă Dumnezeu îmbracă astfel iarba câmpului, care astăzi există, iar mâine este aruncată în cuptor, nu vă va îmbrăca cu mult mai mult pe voi, cei cu puțină credință?
}}


\par }{\PP \VS{31}{\WJ{De aceea, nu vă îngrijorați, zicând: “Ce vom mânca?”, “Ce vom bea?” sau “Cu ce ne vom îmbrăca?”
}}

\VS{32}{\WJ{Căci neamurile caută toate aceste lucruri; dar Tatăl vostru cel ceresc știe că aveți nevoie de toate acestea.
}}

\VS{33}{\WJ{Dar căutați mai întâi Împărăția lui Dumnezeu și dreptatea lui; și toate aceste lucruri vi se vor da și vouă.
}}

\VS{34}{\WJ{De aceea, nu vă îngrijorați pentru ziua de mâine, căci ziua de mâine se va îngrijora singură. Răul propriu al fiecărei zile este suficient.
}}



\par }\Chap{7}{\PP \VerseOne{1}{\WJ{“Nu judecați, ca să nu fiți judecați.
}}

\VS{2}{\WJ{Căci cu orice judecată veți judeca, veți fi judecați; și cu orice măsură veți măsura, vi se va măsura.
}}

\VS{3}{\WJ{De ce vezi paiul care este în ochiul fratelui tău, dar nu te gândești la bârna care este în propriul tău ochi?
}}

\VS{4}{\WJ{Sau cum îi vei spune fratelui tău: “Lasă-mă să scot paiul din ochiul tău”, și iată că bârna este în propriul tău ochi?
}}

\VS{5}{\WJ{Ipocritule! Scoate mai întâi bârna din ochiul tău și apoi vei putea vedea limpede ca să scoți paiul din ochiul fratelui tău.
}}


\par }{\PP \VS{6}{\WJ{“Nu dați la câini ce este sfânt și nu aruncați perlele voastre înaintea porcilor, ca nu cumva să le calce în picioare, să se întoarcă și să vă sfâșie.
}}


\par }{\PP \VS{7}{\WJ{“Cereți și vi se va da. Căutați, și veți găsi. Bateți la ușă și vi se va deschide.
}}

\VS{8}{\WJ{Căci oricine cere primește. Cel care caută găsește. Celui care bate i se va deschide.
}}

\VS{9}{\WJ{Sau cine dintre voi este acela care, dacă fiul său îi cere pâine, îi va da o piatră?
}}

\VS{10}{\WJ{Sau, dacă îi cere un pește, cine îi va da un șarpe?
}}

\VS{11}{\WJ{Deci, dacă voi, care sunteți răi, știți să dați daruri bune copiilor voștri, cu cât mai mult Tatăl vostru, care este în ceruri, va da lucruri bune celor ce i le cer!
}}

\VS{12}{\WJ{Așadar, orice vreți să vă facă vouă oamenii, să le faceți și voi; căci aceasta este Legea și profeții.
}}


\par }{\PP \VS{13}“Intrați
{\WJ{pe poarta cea strâmtă, căci poarta este largă și calea cea largă care duce la pieire, și mulți sunt cei ce intră pe ea.
}}

\VS{14}{\WJ{Cât de strâmtă este poarta și de îngustă este calea care duce la viață! Sunt puțini cei care o găsesc.
}}


\par }{\PP \VS{15}“Feriți-vă
{\WJ{de proorocii mincinoși, care vin la voi în haine de oi, dar pe dinăuntru sunt lupi răpitori.
}}

\VS{16}{\WJ{După roadele lor îi veți cunoaște. Culegeți struguri din spini sau smochine din ciulini?
}}

\VS{17}{\WJ{Tot așa, orice pom bun produce fructe bune, dar pomul stricat produce fructe rele.
}}

\VS{18}{\WJ{Un pom bun nu poate produce fructe rele, și nici un pom corupt nu poate produce fructe bune.
}}

\VS{19}{\WJ{Orice copac care nu produce fructe bune este tăiat și aruncat în foc.
}}

\VS{20}{\WJ{De aceea, după roadele lor îi veți cunoaște.
}}


\par }{\PP \VS{21}“{\WJ{Nu oricine Îmi zice: “Doamne, Doamne!” va intra în Împărăția cerurilor, ci cel ce face voia Tatălui Meu care este în ceruri.
}}

\VS{22}{\WJ{Mulți îmi vor spune în ziua aceea: “Doamne, Doamne, nu cumva am prorocit în numele Tău, nu cumva în numele Tău am scos demoni și nu cumva în numele Tău am făcut multe lucrări puternice?”
}}

\VS{23}{\WJ{Atunci le voi spune: “Nu v-am cunoscut niciodată”. Plecați de la Mine, voi care lucrați fărădelegea'.
}}


\par }{\PP \VS{24}{\WJ{Oricine ascultă aceste cuvinte ale Mele și le împlinește, îl voi asemăna cu un om înțelept care și-a zidit casa pe stâncă.
}}

\VS{25}A căzut
{\WJ{ploaia, au venit inundațiile, au suflat vânturile și au bătut peste casa aceea, dar ea nu a căzut, pentru că era întemeiată pe stâncă.
}}

\VS{26}{\WJ{Oricine aude aceste cuvinte ale mele și nu le pune în practică va fi ca un om nebun care și-a construit casa pe nisip.
}}

\VS{27}{\WJ{Ploaia a căzut, au venit inundațiile, vânturile au suflat și au bătut peste casa aceea; și ea a căzut — și căderea ei a fost mare.”
}}


\par }{\PP \VS{28}După ce a isprăvit Isus de spus acestea, mulțimea a rămas uimită de învățătura Lui,

\VS{29}căci îi învăța cu putere, și nu ca pe cărturari.



\par }\Chap{8}{\PP \VerseOne{1}Când S-a pogorât de pe munte, o mare mulțime de oameni Îl urmau.

\VS{2}Și iată că un lepros a venit la el și i s-a închinat, zicând: “Doamne, dacă vrei, poți să mă faci curat”.


\par }{\PP \VS{3}Isus a întins mâna și l-a atins, zicând: “{\WJ{Vreau să... Să fii curățat”. Imediat, lepra lui a fost curățată.
}}

\VS{4}Isus i-a zis:
{\WJ{“Ai grijă să nu spui nimănui; ci du-te, arată-te preotului și oferă darul pe care l-a poruncit Moise, ca mărturie pentru ei.”
}}


\par }{\PP \VS{5}Când a ajuns în Capernaum, a venit la El un centurion, care i-a cerut ajutor,

\VS{6}zicând: “Doamne, robul meu zace în casă, paralizat și foarte chinuit.”


\par }{\PP \VS{7}Isus i-a zis: “{\WJ{Voi veni și-l voi vindeca.”
}}


\par }{\PP \VS{8}Centurionul a răspuns: “Doamne, nu sunt vrednic să intri sub acoperișul meu. Spune doar un cuvânt și robul meu va fi vindecat.

\VS{9}Căci și eu sunt un om sub autoritate, având sub mine soldați. Îi spun acestuia: “Du-te!” și el se duce; îi spun altuia: “Vino!” și el vine; îi spun robului meu: “Fă aceasta!” și el o face.”


\par }{\PP \VS{10}Isus, auzind aceasta, s-a mirat și a zis celor ce-l urmau: “{\WJ{Adevărat vă spun că n-am găsit o credință atât de mare, nici măcar în Israel.
}}

\VS{11}{\WJ{Vă spun că mulți vor veni de la răsărit și de la apus și vor sta cu Avraam, Isaac și Iacov în Împărăția Cerurilor,
}}

\VS{12}{\WJ{dar copiii Împărăției vor fi aruncați în întunericul de afară. Acolo va fi plâns și scrâșnire de dinți”.”
}}

\VS{13}Isus i-a spus centurionului:
{\WJ{“Du-te și tu. Să ți se facă așa cum ai crezut.” Robul său a fost vindecat în acel ceas.
}}


\par }{\PP \VS{14}Când a intrat Isus în casa lui Petru, a văzut pe mama nevestei lui, care zăcea bolnavă de febră.

\VS{15}I-a atins mâna, și febra a lăsat-o. Ea s-a ridicat și l-a servit.

\VS{16}Când s-a făcut seară, i-au adus pe mulți posedați de demoni. El a izgonit duhurile cu un cuvânt și a vindecat pe toți cei bolnavi,

\VS{17}pentru ca să se împlinească ceea ce s-a spus prin profetul Isaia: “El a luat neputințele noastre și a purtat bolile noastre”.


\par }{\PP \VS{18}Isus, văzând mulțime mare în jurul Lui, a poruncit să se ducă în cealaltă parte.


\par }{\PP \VS{19}Un cărturar a venit și I-a zis: “Învățătorule, Te voi urma oriunde vei merge.”


\par }{\PP \VS{20}Isus i-a zis: “{\WJ{Vulpile au vizuini și păsările cerului au cuiburi, dar Fiul Omului nu are unde să-și pună capul.”
}}


\par }{\PP \VS{21}Un alt ucenic al Lui I-a zis: “Doamne, îngăduie-mi să mă duc mai întâi să îngrop pe tatăl meu.”


\par }{\PP \VS{22}Dar Isus i-a zis:
{\WJ{“Urmează-Mă și lasă morții să-și îngroape morții lor.”
}}


\par }{\PP \VS{23}Când a urcat în corabie, ucenicii Lui L-au urmat.

\VS{24}Iată că s-a iscat o furtună violentă pe mare, încât barca era acoperită de valuri; dar El dormea.

\VS{25}Ucenicii au venit la el și l-au trezit, zicând: “Salvează-ne, Doamne! Suntem pe moarte!”


\par }{\PP \VS{26}El le-a zis: “{\WJ{De ce vă temeți, oameni cu puțină credință?” Apoi s-a ridicat, a mustrat vântul și marea și s-a făcut o mare liniște.
}}


\par }{\PP \VS{27}Și oamenii se mirau și ziceau: “Ce fel de om este acesta, încât și vântul și marea îl ascultă?”


\par }{\PP \VS{28}Când a ajuns de cealaltă parte, în ținutul Gherghesenilor, l-au întâmpinat doi oameni posedați de demoni, care ieșeau din morminte, foarte înverșunați, încât nimeni nu putea trece pe acolo.

\VS{29}Și iată că strigau, zicând: “Ce avem noi de-a face cu Tine, Isus, Fiul lui Dumnezeu? Ai venit aici ca să ne chinuiești înainte de vreme?”

\VS{30}Și era o turmă de mulți porci care păștea departe de ei.

\VS{31}Demonii l-au implorat, spunând: “Dacă ne scoți afară, îngăduie-ne să plecăm în turma de porci.”


\par }{\PP \VS{32}Și le-a zis:
{\WJ{“Mergeți!”
}}

\par }{\PP Au ieșit și au intrat în turma de porci; și iată că întreaga turmă de porci s-a prăbușit în mare și a murit în apă.

\VS{33}Cei care îi hrăneau au fugit și s-au dus în cetate și au povestit totul, inclusiv ce s-a întâmplat cu cei care erau posedați de demoni.

\VS{34}Iată că toată cetatea a ieșit în întâmpinarea lui Isus. Când l-au văzut, l-au implorat să plece de lângă hotarele lor.



\par }\Chap{9}{\PP \VerseOne{1}A intrat într-o corabie, a trecut dincolo și a venit în cetatea sa.

\VS{2}Și iată că i-au adus un om paralizat, întins pe un pat. Isus, văzând credința lor, i-a spus paraliticului:
{\WJ{“Fiule, înveselește-te! Păcatele tale îți sunt iertate”.
}}


\par }{\PP \VS{3}Și iată că unii dintre cărturari își ziceau: “Omul acesta hulește”.


\par }{\PP \VS{4}Isus, cunoscând gândurile lor, a zis: “{\WJ{Pentru ce gândiți rău în inimile voastre?
}}

\VS{5}{\WJ{Căci ce este mai ușor: să spui: “Păcatele tale sunt iertate”, sau să spui: “Scoală-te și umblă?”
}}

\VS{6}{\WJ{Dar ca să știți că Fiul Omului are putere pe pământ să ierte păcatele —
}}(apoi i-a spus paraliticului):
{\WJ{“Scoală-te, ia-ți rogojina și du-te la casa ta”.
}}


\par }{\PP \VS{7}Și s-a sculat și a plecat la casa lui.

\VS{8}Dar mulțimile, văzând aceasta, s-au mirat și au slăvit pe Dumnezeu, care dăduse oamenilor o asemenea putere.


\par }{\PP \VS{9}Când a trecut Isus de acolo, a văzut pe un om numit Matei, care ședea la percepție. I-a spus:
{\WJ{“Urmează-mă”. El s-a ridicat și l-a urmat.
}}

\VS{10}Pe când ședea el în casă, iată că au venit mulți vameși și păcătoși și s-au așezat la masă cu Isus și cu discipolii săi.

\VS{11}Fariseii, văzând aceasta, au zis discipolilor lui: “De ce mănâncă învățătorul vostru cu vameșii și cu păcătoșii?”


\par }{\PP \VS{12}Când a auzit Isus, le-a zis: “{\WJ{Cei sănătoși nu au nevoie de doctor, dar cei bolnavi au nevoie de doctor.
}}

\VS{13}{\WJ{Voi însă mergeți și învățați ce înseamnă acest lucru: “Eu vreau milă, și nu jertfă”, căci n-am venit să chem la pocăință pe cei drepți, ci pe păcătoși.”
}}


\par }{\PP \VS{14}Atunci ucenicii lui Ioan au venit la el și i-au zis: “De ce noi și fariseii postim des, iar ucenicii tăi nu postesc?”


\par }{\PP \VS{15}Isus le-a zis: “{\WJ{Pot oare prietenii mirelui să jelească atâta timp cât mirele este cu ei? Dar vor veni zile când mirele va fi luat de la ei și atunci vor posti.
}}

\VS{16}Nimeni
{\WJ{nu pune o bucată de pânză neîncălzită pe o haină veche, pentru că petecul s-ar rupe de pe haină și se face o gaură și mai mare.
}}

\VS{17}{\WJ{Nici oamenii nu pun vin nou în piei de vin vechi, căci altfel s-ar sparge piei, s-ar vărsa vinul și s-ar strica piei. Nu, ei pun vin nou în piei de vin proaspăt, și amândouă se păstrează.”
}}


\par }{\PP \VS{18}Pe când le spunea acestea, iată că a venit un domnitor și I s-a închinat, zicând: “Fiica mea a murit; vino și pune mâna pe ea și va trăi.”


\par }{\PP \VS{19}Isus s-a sculat și a mers după el, ca și ucenicii lui.

\VS{20}Și iată că o femeie, care avea o scurgere de sânge de doisprezece ani, a venit în urma Lui și s-a atins de marginea hainei Lui;

\VS{21}căci zicea în sinea ei: “Dacă mă ating de haina Lui, mă voi face bine”.


\par }{\PP \VS{22}Dar Isus, întorcându-se și văzând-o, a zis: “{\WJ{Fiică, înveselește-te! Credința ta te-a făcut sănătoasă”. Și femeia s-a făcut bine din acel ceas.
}}


\par }{\PP \VS{23}Când a intrat Isus în casa domnitorului și a văzut pe cei ce cântau la fluier și mulțimea în dezordine,

\VS{24}le-a zis:
{\WJ{“Faceți loc, căci fata nu este moartă, ci doarme.”
}}

\par }{\PP Îl ridiculizau.

\VS{25}Dar, când mulțimea a fost trimisă afară, El a intrat înăuntru, a luat-o de mână și fata s-a ridicat.

\VS{26}Vestea aceasta s-a răspândit în toată țara aceea.


\par }{\PP \VS{27}Pe când trecea Isus de acolo, doi orbi L-au urmărit, strigând și zicând: “Ai milă de noi, Fiul lui David!”

\VS{28}După ce a intrat în casă, orbii s-au apropiat de El. Isus le-a zis:
{\WJ{“Credeți că pot face acest lucru?”
}}

\par }{\PP Ei i-au spus: “Da, Doamne”.


\par }{\PP \VS{29}Apoi s-a atins de ochii lor și a zis: “{\WJ{Fie vouă după credința voastră!”
}}

\VS{30}Atunci li s-au deschis ochii. Isus le-a poruncit cu strictețe, zicând: “{\WJ{Aveți grijă ca nimeni să nu afle despre aceasta.”
}}

\VS{31}Dar ei au ieșit și au răspândit faima Lui în toată țara aceea.


\par }{\PP \VS{32}Pe când ieșeau ei afară, iată că a fost adus la El un mut, care era posedat de un demon.

\VS{33}După ce a fost izgonit demonul, omul mut a vorbit. Mulțimile se minunau și spuneau: “Așa ceva nu s-a mai văzut în Israel!”


\par }{\PP \VS{34}Dar fariseii ziceau: “Prin prințul demonilor scoate demoni.”


\par }{\PP \VS{35}Isus străbătea toate cetățile și satele, învățând în sinagogile lor, propovăduind vestea cea bună a Împărăției și vindecând orice boală și orice neputință în popor.

\VS{36}Când a văzut mulțimile, i s-a făcut milă de ele, pentru că erau hărțuite și împrăștiate, ca niște oi fără păstor.

\VS{37}Atunci le-a spus discipolilor săi: “{\WJ{Într-adevăr, secerișul este mare, dar lucrătorii sunt puțini.
}}

\VS{38}{\WJ{Rugați-vădeci ca Domnul secerișului să trimită lucrători la secerișul Său.”
}}



\par }\Chap{10}{\PP \VerseOne{1}A chemat la Sine pe cei doisprezece ucenici ai Săi și le-a dat putere asupra duhurilor necurate, ca să le scoată afară și să vindece orice boală și orice neputință.

\VS{2}Și numele celor doisprezece apostoli sunt următoarele Cel dintâi, Simon, care se numește Petru; Andrei, fratele său; Iacov, fiul lui Zebedeu; Ioan, fratele său;

\VS{3}Filip; Bartolomeu; Toma; Matei, vameșul; Iacov, fiul lui Alfeu; Lebeu, care se mai numea și Tadeu;

\VS{4}Simon Zelota; și Iuda Iscarioteanul, care l-a și trădat.


\par }{\PP \VS{5}Isus a trimis pe acești doisprezece și le-a poruncit, zicând: “Să
{\WJ{nu mergeți printre neamuri și să nu intrați în nici o cetate a samaritenilor.
}}

\VS{6}{\WJ{Mergeți mai degrabă la oile pierdute ale casei lui Israel.
}}

\VS{7}În
{\WJ{timp ce mergeți, propovăduiți, spunând: “Împărăția Cerurilor este aproape!”
}}

\VS{8}Vindecați
{\WJ{bolnavii, curățați leproșii și alungați demonii. De bunăvoie ați primit, de bunăvoie dați.
}}

\VS{9}{\WJ{Nu luați aur, argint sau aramă în centurile voastre de bani.
}}

\VS{10}{\WJ{Nu luați nici un sac pentru călătorie, nici două haine, nici sandale, nici toiag, căci cel care muncește este vrednic de hrana sa.
}}

\VS{11}{\WJ{În orice cetate sau sat în care intrați, aflați cine este vrednic în el și rămâneți acolo până când veți merge mai departe.
}}

\VS{12}{\WJ{Când intrați în gospodărie, salutați-o.
}}

\VS{13}{\WJ{Dacă gospodăria este vrednică, să vină pacea ta peste ea, iar dacă nu este vrednică, să se întoarcă pacea ta la tine.
}}

\VS{14}{\WJ{Oricine nu vă primește și nu vă ascultă cuvintele, când ieșiți din casa aceea sau din cetatea aceea, scuturați-vă praful de pe picioare.
}}

\VS{15}{\WJ{Adevărat vă spun că, în ziua judecății, va fi mai ușor de suportat pentru țara Sodomei și Gomorei decât pentru orașul acela.
}}


\par }{\PP \VS{16}“{\WJ{Iată, vă trimit ca pe niște oi în mijlocul lupilor. De aceea, fiți înțelepți ca șerpii și inofensivi ca porumbeii.
}}

\VS{17}{\WJ{Dar păziți-vă de oameni, căci vă vor da pe mâna consiliilor și în sinagogile lor vă vor biciui.
}}

\VS{18}{\WJ{Da, și veți fi aduși înaintea guvernatorilor și a împăraților din cauza Mea, ca mărturie pentru ei și pentru neamuri.
}}

\VS{19}{\WJ{Dar, când vă vor preda, nu vă îngrijorați cum și ce veți spune, pentru că vi se va da în ceasul acela ce veți spune.
}}

\VS{20}{\WJ{Căci nu voi vorbiți, ci Duhul Tatălui vostru care vorbește în voi.
}}


\par }{\PP \VS{21}“{\WJ{Fratele va da la moarte pe fratele său și tatăl pe copilul său. Copiii se vor ridica împotriva părinților și îi vor face să moară.
}}

\VS{22}Veți
{\WJ{fi urâți de toți oamenii din cauza numelui Meu, dar cel care va răbda până la sfârșit va fi mântuit.
}}

\VS{23}{\WJ{Dar când vă vor persecuta în această cetate, fugiți în cealaltă, pentru că, cu siguranță, vă spun că nu veți fi trecut prin cetățile lui Israel până nu va veni Fiul Omului.
}}


\par }{\PP \VS{24}“{\WJ{Ucenicul nu este mai presus de învățătorul său, nici sluga mai presus de stăpânul său.
}}

\VS{25}{\WJ{Este de ajuns ca ucenicul să fie ca învățătorul său, și robul ca domnul său. Dacă pe stăpânul casei l-au numit Beelzebul, cu cât mai mult pe cei din casa lui!
}}

\VS{26}{\WJ{De aceea nu vă temeți de ei, căci nu este nimic acoperit care să nu fie descoperit, sau ascuns care să nu fie cunoscut.
}}

\VS{27}{\WJ{Ceea ce vă spun în întuneric, vorbiți la lumină; și ceea ce auziți șoptit la ureche, vestiți pe acoperișurile caselor.
}}

\VS{28}{\WJ{Nu vă temeți de cei care ucid trupul, dar nu pot ucide sufletul. Temeți-vă mai degrabă de cel care este în stare să nimicească și sufletul și trupul în Gheenă.
}}


\par }{\PP \VS{29}{\WJ{“Nu se vând oare două vrăbii cu o monedă de un asarion? Nici una dintre ele nu cade pe pământ fără voia Tatălui vostru.
}}

\VS{30}{\WJ{Dar chiar și firele de păr de pe capul vostru sunt toate numărate.
}}

\VS{31}{\WJ{De aceea nu vă temeți. Voi sunteți mai valoroși decât multe vrăbii.
}}

\VS{32}{\WJ{Așadar, pe oricine mă mărturisește înaintea oamenilor, îl voi mărturisi și eu înaintea Tatălui meu care este în ceruri.
}}

\VS{33}{\WJ{Dar pe oricine Mă va nega înaintea oamenilor, îl voi nega și Eu înaintea Tatălui Meu care este în ceruri.
}}


\par }{\PP \VS{34}{\WJ{“Să nu credeți că am venit să trimit pacea pe pământ. Nu am venit să trimit pace, ci o sabie.
}}

\VS{35}{\WJ{Pentru că am venit să pun pe un om împotriva tatălui său, pe o fiică împotriva mamei sale și pe noră împotriva soacrei sale.
}}

\VS{36}{\WJ{Dușmanii unui om vor fi cei din propria lui casă.
}}

\VS{37}Cine iubește pe
{\WJ{tată sau pe mamă mai mult decât Mine nu este vrednic de Mine; și cine iubește pe fiu sau pe fiică mai mult decât Mine nu este vrednic de Mine.
}}

\VS{38}{\WJ{Cine nu-și ia crucea și nu vine după mine nu este vrednic de mine.
}}

\VS{39}Cine își
{\WJ{caută viața o va pierde, iar cine își pierde viața pentru mine o va găsi.
}}


\par }{\PP \VS{40}{\WJ{Cine vă primește pe voi, pe Mine Mă primește; și cine Mă primește pe Mine, primește pe Cel ce M-a trimis pe Mine.
}}

\VS{41}{\WJ{Cine primește un profet în numele unui profet va primi răsplata unui profet. Cine primește un om drept în numele unui om drept va primi răsplata unui om drept.
}}

\VS{42}{\WJ{Oricineva da unuia dintre acești micuți doar un pahar de apă rece să bea în numele unui ucenic, cu siguranță vă spun că nu-și va pierde nicidecum răsplata.”
}}



\par }\Chap{11}{\PP \VerseOne{1}După ce a isprăvit Isus de a îndruma pe cei doisprezece ucenici ai Săi, a plecat de acolo ca să învețe și să propovăduiască în cetățile lor.


\par }{\PP \VS{2}Ioan, auzind în temniță despre faptele lui Hristos, a trimis doi din ucenicii săi

\VS{3}și I-a zis: “Tu ești Cel ce vine, sau trebuie să căutăm pe altul?”


\par }{\PP \VS{4}Isus le-a răspuns: “{\WJ{Mergeți și spuneți lui Ioan ce ați auzit și ce ați văzut:
}}

\VS{5}{\WJ{orbii își recapătă vederea, șchiopii umblă, leproșii sunt curățiți, surzii aud, morții înviază și săracilor li se vestește o veste bună.
}}

\VS{6}{\WJ{Ferice de cel care nu găsește în mine nici un prilej de poticnire.”
}}


\par }{\PP \VS{7}Pe când se duceau aceștia, Isus a început să spună mulțimii despre Ioan:
{\WJ{“Ce ați ieșit în pustie ca să vedeți? O trestie scuturată de vânt?
}}

\VS{8}{\WJ{Dar voi ce ați ieșit să vedeți? Un om în haine moi? Iată, cei care poartă haine moi sunt în casele regilor.
}}

\VS{9}{\WJ{Dar tu de ce ai ieșit afară? Ca să vezi un profet? Da, vă spun, și chiar mai mult decât un profet.
}}

\VS{10}{\WJ{Căci acesta este Cel despre care este scris: “Iată, trimit înaintea feței tale pe trimisul Meu, care îți va pregăti calea înaintea ta”.
}}

\VS{11}{\WJ{Adevărat vă spun că, dintre cei născuți din femei, nu s-a ridicat nimeni mai mare decât Ioan Botezătorul; totuși, cel mai mic în Împărăția Cerurilor este mai mare decât el.
}}

\VS{12}{\WJ{Din zilele lui Ioan Botezătorul și până acum, Împărăția Cerurilor suferă violențe și cei violenți o iau cu forța.
}}

\VS{13}{\WJ{Căci toți profeții și Legea au profețit până la Ioan.
}}

\VS{14}{\WJ{Dacă sunteți dispuși să-l primiți, acesta este Ilie, care va veni.
}}

\VS{15}{\WJ{Cine are urechi de auzit, să audă.
}}


\par }{\PP \VS{16}{\WJ{Dar cu ce să compar neamul acesta? Seamănă cu niște copii care stau în piețe, care strigă la tovarășii lor
}}

\VS{17}{\WJ{și spun: 'Noi am cântat la flaut pentru voi și voi n-ați dansat. Noi am plâns pentru voi, iar voi nu v-ați plâns'.
}}

\VS{18}{\WJ{Căci Ioan nu a venit nici să mănânce, nici să bea, și ei spun: “Are un demon”.
}}

\VS{19}{\WJ{Fiul Omului a venit mâncând și bând și ei spun: 'Iată un om gurmand și un bețiv, prietenul vameșilor și al păcătoșilor!'. Dar înțelepciunea este justificată de copiii ei.”
}}


\par }{\PP \VS{20}Apoi a început să denunțe cetățile în care se săvârșiseră cele mai multe din faptele lui mărețe, pentru că nu se pocăiseră.

\VS{21}{\WJ{“Vai de tine, Corazin! Vai de tine, Betsaida! Pentru că, dacă în Tir și în Sidon s-ar fi făcut în Tir și în Sidon lucrările mărețe care s-au făcut în voi, s-ar fi pocăit demult în sac și cenușă.
}}

\VS{22}{\WJ{Dar vă spun că, în ziua judecății, va fi mai ușor de suportat pentru Tir și Sidon decât pentru voi.
}}

\VS{23}{\WJ{Tu, Capernaum, care te-ai înălțat la cer, te vei coborî în Hades. Căci, dacă s-ar fi făcut în Sodoma lucrările mărețe care s-au făcut în voi, ar fi rămas până astăzi.
}}

\VS{24}{\WJ{Dar Eu vă spun că, în ziua judecății, va fi mai ușor de suportat pentru țara Sodomei decât pentru voi.”
}}


\par }{\PP \VS{25}În vremea aceea, Isus a răspuns: “Îți
{\WJ{mulțumesc, Tată, Doamne al cerului și al pământului, că ai ascuns aceste lucruri de cei înțelepți și pricepuți și le-ai descoperit copiilor.
}}

\VS{26}{\WJ{Da, Tată, pentru că așa a fost bineplăcut înaintea Ta.
}}

\VS{27}{\WJ{Toate lucrurile mi-au fost predate de Tatăl meu. Nimeni nu-L cunoaște pe Fiul, în afară de Tatăl; și nimeni nu-L cunoaște pe Tatăl, în afară de Fiul și de cel căruia Fiul vrea să i-L dezvăluie.
}}


\par }{\PP \VS{28}“{\WJ{Veniți la Mine, voi toți cei osteniți și împovărați, și Eu vă voi da odihnă.
}}

\VS{29}Luați
{\WJ{jugul Meu asupra voastră și învățați de la Mine, căci Eu sunt blând și smerit cu inima; și veți găsi odihnă pentru sufletele voastre.
}}

\VS{30}{\WJ{Căcijugul Meu este ușor și povara Mea este ușoară.”
}}



\par }\Chap{12}{\PP \VerseOne{1}În vremea aceea, Iisus mergea în ziua de Sabat prin lanurile de grâu. Discipolilor Săi le era foame și au început să culeagă căpățâni de grâu și să mănânce.

\VS{2}Dar fariseii, văzând aceasta, i-au spus: “Iată că discipolii tăi fac ceea ce nu este îngăduit să facă în ziua de Sabat.”


\par }{\PP \VS{3}Dar El le-a zis: “{\WJ{N-ați citit ce a făcut David și cei ce erau cu el, când i-a fost foame?
}}

\VS{4}{\WJ{cum a intrat în casa lui Dumnezeu și a mâncat pâinea de spectacol, pe care nu era îngăduit să o mănânce nici el, nici cei care erau cu el, ci numai preoții?
}}

\VS{5}{\WJ{Sau n-ați citit în Lege că, în ziua Sabatului, preoții din templu profanează Sabatul și sunt nevinovați?
}}

\VS{6}{\WJ{Dar Eu vă spun că unul mai mare decât templul este aici.
}}

\VS{7}{\WJ{Dar dacă ați fi știut ce înseamnă acest lucru: “Eu doresc milă, și nu jertfă”, nu i-ați fi condamnat pe cei fără vină.
}}

\VS{8}{\WJ{Căci Fiul Omului este Domnul Sabatului.”
}}


\par }{\PP \VS{9}A plecat de acolo și a intrat în sinagoga lor.

\VS{10}Și iată că era acolo un om cu o mână uscată. L-au întrebat: “Este îngăduit să vindeci în ziua de Sabat?”, ca să-l acuze.


\par }{\PP \VS{11}El le-a zis: “Cine
{\WJ{dintre voi are o oaie și, dacă aceasta cade într-o groapă în ziua de Sabat, nu se va apuca de ea și n-o va scoate afară?
}}

\VS{12}{\WJ{Cu cât este mai valoros atunci un om decât o oaie! De aceea este permis să faci binele în ziua de Sabat”.
}}

\VS{13}Apoi i-a spus omului:
{\WJ{“Întinde-ți mâna”. El a întins-o; și ea a fost refăcută întreagă, la fel ca cealaltă.
}}

\VS{14}Dar fariseii au ieșit și au uneltit împotriva lui, ca să-l nimicească.


\par }{\PP \VS{15}Isus, aflând aceasta, s-a retras de acolo. L-au urmat mulțimi mari de oameni; și El i-a vindecat pe toți,

\VS{16}și le-a poruncit să nu-L facă cunoscut,

\VS{17}ca să se împlinească ceea ce fusese spus prin profetul Isaia, care zicea


\par }{\Q \VS{18}“Iată robul Meu, pe care l-am ales,

\par }{\QB iubitul meu, în care sufletul meu este binevoitor.

\par }{\Q Voi pune Duhul Meu peste el.

\par }{\QB El va proclama dreptatea națiunilor.


\par }{\Q \VS{19}El nu se va zbate și nu va striga,

\par }{\QB nici nu-i va auzi cineva vocea pe străzi.


\par }{\Q \VS{20}El nu va rupe trestia zdrobită.

\par }{\QB El nu va stinge nici un fân care fumegă,

\par }{\Q până când va conduce justiția la victorie.


\par }{\QB \VS{21}În numele Lui vor nădăjdui neamurile.”


\par }{\PP \VS{22}Atunci a fost adus la El un orb și mut, posedat de un demon, și L-a vindecat, astfel că orbul și mutul vorbea și vedea.

\VS{23}Toată mulțimea era uimită și zicea: “Oare poate fi acesta fiul lui David?”

\VS{24}Dar fariseii, când au auzit, au zis: “Omul acesta nu scoate demoni decât prin Beelzebul, prințul demonilor.”


\par }{\PP \VS{25}Cunoscându-le gândurile, Isus le-a zis: “{\WJ{Orice împărăție împărțită împotriva ei însăși este dărâmată, și orice cetate sau casă împărțită împotriva ei însăși nu va rămâne în picioare.
}}

\VS{26}{\WJ{Dacă Satana îl alungă pe Satana, el este dezbinat împotriva lui însuși. Cum va rămâne atunci împărăția lui în picioare?
}}

\VS{27}{\WJ{Dacă Eu, prin Beelzebul, scot afară demoni, prin cine îi scot copiii voștri? De aceea ei vor fi judecătorii voștri.
}}

\VS{28}{\WJ{Dar dacă eu, prin Duhul lui Dumnezeu, scot afară demoni, atunci Împărăția lui Dumnezeu a venit peste voi.
}}

\VS{29}{\WJ{Sau cum poate cineva să intre în casa celui puternic și să-i jefuiască bunurile, dacă nu-l leagă mai întâi pe cel puternic? Atunci îi va jefui casa.
}}


\par }{\PP \VS{30}Cine
{\WJ{nu este cu Mine, este împotriva Mea, și cine nu se adună cu Mine, risipește.
}}

\VS{31}{\WJ{De aceea vă spun că orice păcat și orice blasfemie va fi iertat oamenilor, dar blasfemia împotriva Duhului nu va fi iertată oamenilor.
}}

\VS{32}{\WJ{Oricine va spune un cuvânt împotriva Fiului Omului, îi va fi iertat, dar oricine va vorbi împotriva Duhului Sfânt, nu-i va fi iertat, nici în veacul acesta, nici în cel viitor.
}}


\par }{\PP \VS{33}“Ori faceți pomul
{\WJ{bun și rodul lui bun, ori faceți pomul stricat și rodul lui stricat, căci pomul se cunoaște după rodul lui.
}}

\VS{34}{\WJ{Voi, urmași ai viperelor, cum puteți, fiind răi, să vorbiți lucruri bune? Căci din belșugul inimii, gura vorbește.
}}

\VS{35}{\WJ{Omul bun din comoara sa bună scoate lucruri bune, iar omul rău din comoara sa rea scoate lucruri rele.
}}

\VS{36}{\WJ{Eu vă spun că orice cuvânt deșert pe care-l spun oamenii, vor da socoteală de el în ziua judecății.
}}

\VS{37}{\WJ{Căci prin cuvintele voastre veți fi îndreptățiți, iar prin cuvintele voastre veți fi osândiți.”
}}


\par }{\PP \VS{38}Atunci unii dintre cărturari și farisei au răspuns: “Învățătorule, vrem să vedem un semn de la Tine.”


\par }{\PP \VS{39}Dar El le-a răspuns: “{\WJ{Un neam rău și adulterin caută un semn; dar nu i se va da alt semn decât semnul proorocului Iona.
}}

\VS{40}{\WJ{Căci, după cum Iona a stat trei zile și trei nopți în pântecele peștelui uriaș, tot așa și Fiul Omului va sta trei zile și trei nopți în inima pământului.
}}

\VS{41}{\WJ{Bărbații din Ninive vor sta în picioare la judecată cu acest neam și îl vor condamna, pentru că s-au pocăit la propovăduirea lui Iona; și iată că cineva mai mare decât Iona este aici.
}}

\VS{42}{\WJ{Împărăteasa Sudului se va ridica la judecata cu acest neam și-l va osândi, pentru că a venit de la marginile pământului să asculte înțelepciunea lui Solomon; și iată, cineva mai mare decât Solomon este aici.
}}


\par }{\PP \VS{43}{\WJ{“Când un duh necurat iese din om, trece prin locuri fără apă, căutând odihnă, și nu o găsește.
}}

\VS{44}{\WJ{Atunci el spune: “Mă voi întoarce în casa mea de unde am venit.” Și când se întoarce, o găsește goală, măturată și pusă în ordine.
}}

\VS{45}{\WJ{Atunci se duce și ia cu el alte șapte duhuri mai rele decât el, intră înăuntru și locuiesc acolo. Ultima stare a acelui om devine mai rea decât prima. Tot așa va fi și cu această generație rea.”
}}


\par }{\PP \VS{46}Pe când vorbea El încă mulțimilor, iată că mama și frații Lui stăteau afară și căutau să vorbească cu El.

\VS{47}Cineva i-a zis: “Iată, mama ta și frații tăi stau afară și caută să vorbească cu tine.”


\par }{\PP \VS{48}Dar el a răspuns celui ce-i vorbea:
{\WJ{“Cine este mama mea? Cine sunt frații mei?”
}}

\VS{49}El a întins mâna spre ucenicii Săi și a zis:
{\WJ{“Iată mama Mea și frații Mei!
}}

\VS{50}{\WJ{Căci oricine face voia Tatălui Meu, care este în ceruri, acela este fratele Meu, sora Mea și mama Mea.”
}}



\par }\Chap{13}{\PP \VerseOne{1}În ziua aceea, Isus a ieșit din casă și a șezut pe malul mării.

\VS{2}O mare mulțime s-a adunat la El, așa că a intrat într-o barcă și a șezut; și toată mulțimea stătea pe plajă.

\VS{3}El le vorbea multe lucruri în parabole, zicând:
{\WJ{“Iată, un agricultor a ieșit să semene.
}}

\VS{4}În
{\WJ{timp ce semăna, unele semințe au căzut pe marginea drumului, iar păsările au venit și le-au mâncat.
}}

\VS{5}{\WJ{Altele au căzut pe un teren stâncos, unde nu prea avea pământ, și imediat au răsărit, pentru că nu aveau pământ adânc.
}}

\VS{6}{\WJ{Când a răsărit soarele, s-au pârjolit. Pentru că nu aveau rădăcini, s-au uscat.
}}

\VS{7}{\WJ{Altele au căzut între spini. Spinii au crescut și le-au sufocat.
}}

\VS{8}{\WJ{Altele au căzut pe pământ bun și au dat rod: unele de o sută de ori mai mult, altele de șaizeci, altele de treizeci.
}}

\VS{9}{\WJ{Cine are urechi de auzit, să asculte.”
}}


\par }{\PP \VS{10}Ucenicii au venit și I-au zis: “De ce le vorbești în pilde?”


\par }{\PP \VS{11}El le-a răspuns: “{\WJ{Vouă v-a fost dat să cunoașteți tainele Împărăției cerurilor, dar lor nu le este dat.
}}

\VS{12}{\WJ{Căci oricui are, i se va da și va avea din belșug; dar oricui nu are, i se va lua chiar și ceea ce are.
}}

\VS{13}{\WJ{De aceea le vorbesc în parabole, pentru că, văzând, nu văd, iar auzind, nu aud și nici nu înțeleg.
}}

\VS{14}{\WJ{În ei se împlinește profeția lui Isaia, care zice:}},

\par }{\Q {\WJ{'Auzind, veți auzi,
}}

\par }{\QB {\WJ{și nu va înțelege în niciun fel;
}}

\par }{\Q {\WJ{Văzând vei vedea,
}}

\par }{\QB {\WJ{și nu va percepe în niciun fel;
}}


\par }{\Q \VS{15}{\WJ{Căci inima acestui popor a devenit insensibilă,
}}

\par }{\QB {\WJ{urechile lor sunt surde de auz,
}}

\par }{\QB {\WJ{și au închis ochii;
}}

\par }{\Q {\WJ{sau poate că ar putea percepe cu ochii,
}}

\par }{\QB {\WJ{aud cu urechile lor,
}}

\par }{\QB {\WJ{să înțeleagă cu inima lor,
}}

\par }{\Q {\WJ{și s-ar întoarce din nou,
}}

\par }{\QB {\WJ{și eu îi voi vindeca.
}}


\par }{\PP \VS{16}“{\WJ{Binecuvântați sunt ochii voștri, căci văd, și urechile voastre, căci aud.
}}

\VS{17}{\WJ{Căci cu adevărat vă spun că mulți profeți și oameni drepți au dorit să vadă lucrurile pe care le vedeți voi și nu le-au văzut; și să audă lucrurile pe care le auziți voi și nu le-au auzit.
}}


\par }{\PP \VS{18}“{\WJ{Ascultați, deci, pilda țăranului.
}}

\VS{19}{\WJ{Când cineva aude cuvântul Împărăției și nu-l înțelege, vine cel rău și smulge ceea ce a fost semănat în inima lui. Iată ce a fost semănat la marginea drumului.
}}

\VS{20}{\WJ{Ceea ce a fost semănat pe locurile stâncoase, acesta este cel care aude cuvântul și îndată îl primește cu bucurie;
}}

\VS{21}{\WJ{dar nu are rădăcină în el însuși, ci rezistă o vreme. Când apare opresiunea sau persecuția din cauza cuvântului, îndată se poticnește.
}}

\VS{22}{\WJ{Ceea ce a fost semănat între spini, acesta este cel care ascultă cuvântul, dar grijile veacului acestuia și înșelăciunea bogățiilor înăbușă cuvântul, și el devine neînsuflețit.
}}

\VS{23}{\WJ{Ceea ce a fost semănat în pământ bun, acesta este cel care aude cuvântul și îl înțelege, care cu siguranță face rod și produce, unii de o sută de ori mai mult, alții de șaizeci, iar alții de treizeci.”
}}


\par }{\PP \VS{24}Și le-a pus înaintea lor o altă pildă, zicând: “{\WJ{Împărăția cerurilor se aseamănă cu un om care a semănat sămânță bună în câmpul său;
}}

\VS{25}{\WJ{dar, pe când dormeau oamenii, a venit vrăjmașul său, a semănat și el buruieni de neghină printre grâu și a plecat.
}}

\VS{26}{\WJ{Dar când a răsărit firul și a produs grâu, atunci au apărut și buruienile de scorbură.
}}

\VS{27}{\WJ{Slujitorii gospodarului au venit și i-au spus: “Domnule, nu ai semănat tu sămânță bună în câmpul tău? De unde au apărut aceste buruieni de scorbură?”.
}}


\par }{\PP \VS{28}{\WJ{El le-a zis: “Un vrăjmaș a făcut aceasta.
}}

\par }{\PP {\WJ{Slujitorii l-au întrebat: “Vrei să mergem noi să-i adunăm?}}”.


\par }{\PP \VS{29}{\WJ{Dar El a zis: “Nu, ca nu cumva, adunând buruienile, să nu dezrădăcinați cu ele grâul.
}}

\VS{30}{\WJ{Lasă-le pe amândouă să crească împreună până la seceriș, iar la vremea secerișului voi spune secerătorilor: “Mai întâi, strângeți buruienile de scorbură și legați-le în mănunchiuri ca să le ardeți; dar grâul strângeți-l în hambarul Meu."”””.
}}


\par }{\PP \VS{31}Și le-a pus înainte o altă pildă, zicând: “{\WJ{Împărăția cerurilor este ca un grăunte de muștar pe care l-a luat un om și l-a semănat în câmpul său,
}}

\VS{32}{\WJ{care este mai mic decât toate semințele. Dar, când a crescut, este mai mare decât ierburile și devine un copac, astfel încât păsările cerului vin și se adăpostesc în ramurile lui.”
}}


\par }{\PP \VS{33}Și le-a spus o altă pildă.
{\WJ{“Împărăția cerurilor se aseamănă cu drojdia pe care a luat-o o femeie și a ascuns-o în trei măsuri de făină, până ce a fost toată leșinată.”
}}


\par }{\PP \VS{34}Isus vorbea mulțimilor toate aceste lucruri în pilde; și fără pilde nu le vorbea,

\VS{35}ca să se împlinească ce fusese spus prin proorocul care zicea,

\par }{\Q “Îmi voi deschide gura în parabole;

\par }{\QB Voi descoperi lucruri ascunse de la întemeierea lumii.”


\par }{\PP \VS{36}Atunci Isus a dat drumul mulțimii și a intrat în casă. Ucenicii Lui au venit la El și I-au zis: “Explică-ne parabola buruienilor de pe câmp.”


\par }{\PP \VS{37}El le-a răspuns: “{\WJ{Cel ce seamănă sămânța cea bună este Fiul Omului;
}}

\VS{38}{\WJ{câmpul este lumea, semințele bune sunt copiii Împărăției, iar buruienile sunt copiii celui rău.
}}

\VS{39}{\WJ{Dușmanul care le-a semănat este diavolul. Secerișul este sfârșitul veacului, iar secerătorii sunt îngerii.
}}

\VS{40}De
{\WJ{aceea, așa cum buruienile de scorbură sunt adunate și arse în foc, așa va fi și la sfârșitul acestui veac.
}}

\VS{41}{\WJ{Fiul Omului va trimite pe îngerii Săi, care vor strânge din Împărăția Sa tot ce este de poticnire și pe cei ce fac nelegiuiri,
}}

\VS{42}{\WJ{și îi vor arunca în cuptorul de foc. Acolo va fi plâns și scrâșnire de dinți.
}}

\VS{43}{\WJ{Atunci cei drepți vor străluci ca soarele în Împărăția Tatălui lor. Cine are urechi de auzit, să audă.
}}


\par }{\PP \VS{44}{\WJ{“Împărăția cerurilor se aseamănă cu o comoară ascunsă în câmp, pe care a găsit-o un om și a ascuns-o. În bucuria lui, se duce și vinde tot ce are și cumpără acel câmp.
}}


\par }{\PP \VS{45}{\WJ{“Împărăția cerurilor se aseamănă cu un om care este un negustor care caută perle frumoase,
}}

\VS{46}{\WJ{și care, găsind o perlă de mare preț, s-a dus, a vândut tot ce avea și a cumpărat-o.
}}


\par }{\PP \VS{47}{\WJ{“Împărăția cerurilor se aseamănă cu o năvodă aruncată în mare, care a adunat pești de toate felurile,
}}

\VS{48}și pe care,
{\WJ{când s-a umplut, pescarii au tras-o pe plajă. S-au așezat și au adunat ce era bun în recipiente, dar ce era rău l-au aruncat.
}}

\VS{49}{\WJ{Așa va fi și la sfârșitul lumii. Îngerii vor veni și îi vor despărți pe cei răi dintre cei drepți
}}

\VS{50}{\WJ{și îi vor arunca în cuptorul de foc. Acolo va fi plâns și scrâșnire de dinți.”
}}

\VS{51}Isus le-a zis:
{\WJ{“Ați înțeles toate aceste lucruri?”
}}

\par }{\PP Ei i-au răspuns: “Da, Doamne”.


\par }{\PP \VS{52}El le-a zis: “{\WJ{De aceea, orice cărturar care a fost făcut ucenic în Împărăția cerurilor se aseamănă cu un om gospodar, care scoate din comoara sa lucruri noi și vechi.”
}}


\par }{\PP \VS{53}După ce a isprăvit Isus aceste pilde, a plecat de acolo.

\VS{54}Venind în țara lui, i-a învățat pe oameni în sinagoga lor, așa că ei erau uimiți și ziceau: “De unde are omul acesta atâta înțelepciune și atâtea minuni?

\VS{55}Nu cumva este acesta fiul tâmplarului? Nu cumva mama lui se numește Maria, iar frații săi Iacov, Iose, Simon și Iuda?

\VS{56}Nu sunt toate surorile lui cu noi? De unde a obținut atunci acest om toate aceste lucruri?”.

\VS{57}Ei s-au simțit jigniți de el.

\par }{\PP Dar Isus le-a zis: “{\WJ{Un profet nu este lipsit de cinste decât în țara și în casa lui.”
}}

\VS{58}El nu a făcut acolo multe lucrări mărețe din cauza necredinței lor.



\par }\Chap{14}{\PP \VerseOne{1}În vremea aceea, Irod tetrarhul a auzit vestea despre Isus

\VS{2}și a zis slujitorilor săi: “Acesta este Ioan Botezătorul. El a înviat din morți. De aceea lucrează în el aceste puteri”.

\VS{3}Pentru că Irod îl prinsese pe Ioan, îl legase și îl întemnițase din cauza Irodiadei, soția fratelui său Filip.

\VS{4}Căci Ioan îi spusese: “Nu-ți este îngăduit să o ai.”

\VS{5}Când a vrut să-l omoare, s-a temut de mulțime, pentru că îl socotea profet.

\VS{6}Dar când a venit ziua de naștere a lui Irod, fiica Irodiadei a dansat printre ei și i-a plăcut lui Irod.

\VS{7}De aceea a promis cu jurământ că-i va da tot ce-i va cere.

\VS{8}Ea, fiind îndemnată de mama ei, a zis: “Dă-mi aici, pe un platou, capul lui Ioan Botezătorul”.


\par }{\PP \VS{9}Împăratul s-a mâhnit, dar a poruncit să fie dat, pentru jurămintele sale și pentru cei ce stăteau la masă cu el.

\VS{10}A trimis și a decapitat pe Ioan în temniță.

\VS{11}Capul lui a fost adus pe un platou și dat tinerei domnișoare, iar ea l-a dus mamei sale.

\VS{12}Ucenicii lui au venit, au luat trupul și l-au îngropat. Apoi s-au dus și au anunțat pe Isus.

\VS{13}Isus, când a auzit aceasta, s-a retras de acolo cu o barcă într-un loc pustiu, la o parte. Când au auzit mulțimile, l-au urmat pe jos dinspre cetăți.


\par }{\PP \VS{14}Isus a ieșit afară și a văzut o mulțime mare. I s-a făcut milă de ei și le-a vindecat bolnavii.

\VS{15}Când s-a făcut seară, ucenicii Lui au venit la El și I-au zis: “Locul acesta este pustiu și ora este deja târzie. Trimiteți mulțimea să plece, ca să se ducă în sate și să-și cumpere de mâncare.”


\par }{\PP \VS{16}Dar Isus le-a zis: “Nu este
{\WJ{nevoie să plece. Dați-le voi ceva de mâncare”.
}}


\par }{\PP \VS{17}Ei I-au spus: “Avem aici numai cinci pâini și doi pești.”


\par }{\PP \VS{18}El a zis: “{\WJ{Adu-i aici la Mine”.
}}

\VS{19}A poruncit mulțimilor să se așeze pe iarbă; apoi a luat cele cinci pâini și cei doi pești și, privind spre cer, a binecuvântat, a frânt și a dat pâinile ucenicilor, iar ucenicii au dat mulțimilor.

\VS{20}Toți au mâncat și s-au săturat. Au luat douăsprezece coșuri pline cu ceea ce rămăsese de la bucățile sparte.

\VS{21}Cei care au mâncat erau cam cinci mii de bărbați, pe lângă femei și copii.


\par }{\PP \VS{22}Îndată Isus a pus pe ucenici să se urce în corabie și să meargă înaintea Lui, în cealaltă parte, în timp ce El gonea mulțimea.

\VS{23}După ce a alungat mulțimile, s-a urcat singur pe munte ca să se roage. Când a venit seara, era acolo singur.

\VS{24}Dar corabia se afla acum în mijlocul mării, tulburată de valuri, căci vântul era împotrivă.

\VS{25}În a patra veghe a nopții, Isus a venit la ei, mergând pe mare.

\VS{26}Când discipolii l-au văzut umblând pe mare, s-au tulburat și au zis: “Este o fantomă!” și au strigat de frică.

\VS{27}Dar îndată Isus le-a vorbit și le-a zis:
{\WJ{“Înveseliți-vă! Eu sunt! Nu vă temeți!”.
}}


\par }{\PP \VS{28}Petru I-a răspuns: “Doamne, dacă ești Tu, poruncește-mi să vin la Tine pe ape.”


\par }{\PP \VS{29}El a zis:
{\WJ{“Vino!”
}}

\par }{\PP Petru a coborât din barcă și a mers pe ape pentru a veni la Isus.

\VS{30}Dar, văzând că vântul era puternic, s-a temut și, începând să se scufunde, a strigat: “Doamne, salvează-mă!”.


\par }{\PP \VS{31}Îndată Isus a întins mâna, l-a apucat și i-a zis: “{\WJ{Tu, cel cu puțină credință, de ce te-ai îndoit?”
}}

\VS{32}Când s-au urcat în barcă, vântul a încetat.

\VS{33}Cei care erau în barcă au venit și i s-au închinat, spunând: “Cu adevărat ești Fiul lui Dumnezeu!”


\par }{\PP \VS{34}După ce au trecut dincolo, au ajuns în ținutul Ghenaretului.

\VS{35}Când oamenii din acel loc l-au recunoscut, au trimis în toată regiunea din jur și i-au adus pe toți cei bolnavi;

\VS{36}și l-au rugat să atingă doar franjurii hainei lui. Toți cei care o atingeau se vindecau.



\par }\Chap{15}{\PP \VerseOne{1}Fariseii și cărturarii au venit de la Ierusalim la Isus, și au zis:

\VS{2}“Pentru ce nu ascultă ucenicii Tăi tradiția bătrânilor? Căci ei nu se spală pe mâini când mănâncă pâine”.


\par }{\PP \VS{3}El le-a răspuns: “Pentru
{\WJ{ce nu ascultați și voi de porunca lui Dumnezeu, din pricina tradiției voastre?
}}

\VS{4}{\WJ{Căci Dumnezeu a poruncit: “Cinstește pe tatăl tău și pe mama ta” și: “Cine vorbește de rău pe tatăl sau pe mama, să fie omorât”.
}}

\VS{5}{\WJ{Dar voi ziceți: “Oricine ar putea spune tatălui său sau mamei sale: “Orice ajutor pe care l-ai fi putut primi altfel de la mine este un dar închinat lui Dumnezeu”,
}}

\VS{6}{\WJ{să nu-și cinstească tatăl sau mama”. Voi ați făcut ca porunca lui Dumnezeu să fie nulă din cauza tradiției voastre.
}}

\VS{7}{\WJ{Voi, ipocriților! Bine a profețit Isaia despre voi, zicând: “Nu vă faceți de râs!
}}


\par }{\Q \VS{8}{\WJ{“Oamenii aceștia se apropie de Mine cu gura lor,
}}

\par }{\QB {\WJ{și mă onorează cu buzele lor;
}}

\par }{\QB {\WJ{dar inima lor este departe de mine.
}}


\par }{\Q \VS{9}Și
{\WJ{în zadar se închină la Mine,
}}

\par }{\QB {\WJ{învățând ca doctrină reguli făcute de oameni.””
}}


\par }{\PP \VS{10}A chemat mulțimea și le-a zis: “{\WJ{Ascultați și înțelegeți.
}}

\VS{11}Ceea ce
{\WJ{intră în gură nu spurcă pe om; dar ceea ce iese din gură, aceasta spurcă pe om.”
}}


\par }{\PP \VS{12}Atunci au venit ucenicii și I-au zis: “Știi că fariseii s-au mâniat când au auzit cuvântul acesta?”


\par }{\PP \VS{13}Dar El a răspuns: “{\WJ{Orice plantă pe care nu a plantat-o Tatăl Meu cel ceresc va fi dezrădăcinată.
}}

\VS{14}{\WJ{Lăsați-le în pace. Ei sunt călăuze oarbe ale orbilor. Dacă orbul îi călăuzește pe orbi, amândoi vor cădea într-o groapă.”
}}


\par }{\PP \VS{15}Petru i-a răspuns: “Explică-ne pilda aceasta.”


\par }{\PP \VS{16}Și Isus a zis: “{\WJ{Nu înțelegeți nici voi încă?
}}

\VS{17}{\WJ{Nu înțelegeți că tot ce intră în gură trece în pântece și apoi iese din trup?
}}

\VS{18}{\WJ{Dar ceea ce iese din gură iese din inimă și îl spurcă pe om.
}}

\VS{19}{\WJ{Căci din inimă ies gândurile rele, uciderile, adulterele, păcatele sexuale, furturile, mărturia mincinoasă și blasfemiile.
}}

\VS{20}{\WJ{Acestea sunt lucrurile care îl spurcă pe om; dar a mânca cu mâinile nespălate nu-l spurcă pe om.”
}}


\par }{\PP \VS{21}Isus a ieșit de acolo și s-a retras în ținutul Tirului și al Sidonului.

\VS{22}Și iată că o femeie cananeeancă a ieșit din acele hotare și a strigat: “Ai milă de mine, Doamne, Fiul lui David! Fiica mea este grav posedată de un demon!”


\par }{\PP \VS{23}Dar el nu i-a răspuns nici un cuvânt.

\par }{\PP Ucenicii Lui au venit și L-au rugat, zicând: “Trimite-o departe, căci strigă după noi”.


\par }{\PP \VS{24}Dar El a răspuns: “{\WJ{Eu n-am fost trimis decât la oile pierdute ale casei lui Israel.”
}}


\par }{\PP \VS{25}Dar ea a venit și s-a închinat Lui, zicând: “Doamne, ajută-mă!”


\par }{\PP \VS{26}Iar el a răspuns: “{\WJ{Nu se cuvine să iei pâinea copiilor și să o arunci la câini.”
}}


\par }{\PP \VS{27}Iar ea a răspuns: “Da, Doamne, dar și câinii mănâncă firimiturile care cad de la masa stăpânilor lor.”


\par }{\PP \VS{28}Atunci Isus i-a răspuns: “{\WJ{Femeie, mare este credința ta! Să ți se facă așa cum dorești!” Și fiica ei a fost vindecată din acel ceas.
}}


\par }{\PP \VS{29}Isus a plecat de acolo și s-a apropiat de marea Galileii; apoi S-a suit pe munte și a șezut acolo.

\VS{30}Mulțimi mari au venit la El, având cu ei șchiopi, orbi, muți, mutilați și mulți alții, și i-au pus la picioare. El i-a vindecat,

\VS{31}astfel încât mulțimea se mira când vedea pe cei muți vorbind, pe cei răniți vindecându-se, pe cei șchiopi mergând și pe cei orbi văzând — și slăveau pe Dumnezeul lui Israel.


\par }{\PP \VS{32}Isus a chemat pe ucenicii Săi și le-a zis: “Mi-e
{\WJ{milă de mulțime, pentru că de trei zile stau cu Mine și nu au ce să mănânce. Nu vreau să îi trimit departe în post, pentru că ar putea să leșine pe drum.”
}}


\par }{\PP \VS{33}Ucenicii I-au zis: “De unde am putea lua atâtea pâini într-un loc pustiu ca să săturăm o mulțime atât de mare?”


\par }{\PP \VS{34}Isus le-a zis:
{\WJ{“Câte pâini aveți?”
}}

\par }{\PP Ei au spus: “Șapte și câțiva peștișori”.


\par }{\PP \VS{35}Și a poruncit mulțimii să șadă pe pământ;

\VS{36}și a luat cele șapte pâini și peștii. A mulțumit și le-a frânt; le-a împărțit și le-a dat ucenicilor, iar ucenicii mulțimii.

\VS{37}Toți au mâncat și s-au săturat. Au luat șapte coșuri pline cu bucățile sparte care rămăseseră.

\VS{38}Cei care au mâncat erau patru mii de bărbați, pe lângă femei și copii.

\VS{39}Apoi a alungat mulțimile, s-a urcat în barcă și a ajuns în hotarele Magdala.



\par }\Chap{16}{\PP \VerseOne{1}Fariseii și saducheii au venit și, punându-L la încercare, L-au rugat să le arate un semn din cer.

\VS{2}Dar el le-a răspuns: “{\WJ{Când se face seară, voi spuneți: “Va fi vreme frumoasă, căci cerul este roșu”.
}}

\VS{3}{\WJ{Dimineața: 'Astăzi va fi vreme rea, căci cerul este roșu și amenințător'. Ipocriților! Știți să discerneți aspectul cerului, dar nu puteți discerne semnele vremurilor!
}}

\VS{4}{\WJ{O generație rea și adulterină caută un semn, dar nu i se va da niciun semn, cu excepția semnului profetului Iona.”
}}

\par }{\PP I-a lăsat și a plecat.

\VS{5}Ucenicii au ajuns de partea cealaltă și au uitat să ia pâine.

\VS{6}Isus le-a zis: “Luați
{\WJ{aminte și feriți-vă de drojdia fariseilor și a saducheilor.”
}}


\par }{\PP \VS{7}Și se gândeau între ei, zicând: “Nu am adus pâine.”


\par }{\PP \VS{8}Isus, văzând aceasta, a zis: “{\WJ{De ce vă certați între voi, oameni cu puțină credință, pentru că n-ați adus pâine?
}}

\VS{9}{\WJ{Încă nu vă dați seama și nu vă amintiți de cele cinci pâini pentru cele cinci mii și de câte coșuri ați luat,
}}

\VS{10}{\WJ{sau de cele șapte pâini pentru cele patru mii și de câte coșuri ați luat?
}}

\VS{11}{\WJ{Cum se face că nu vă dați seama că nu v-am vorbit despre pâini? Ci feriți-vă de drojdia fariseilor și a saducheilor.”
}}


\par }{\PP \VS{12}Atunci au înțeles că nu le spusese să se ferească de drojdia pâinii, ci de învățătura fariseilor și a saducheilor.


\par }{\PP \VS{13}Când a ajuns Isus în părțile Cezareei lui Filipi, a întrebat pe ucenicii Săi, zicând: “{\WJ{Cine zic oamenii că sunt Eu, Fiul Omului?”
}}


\par }{\PP \VS{14}Ei au zis: “Unii zic Ioan Botezătorul, alții Ilie, iar alții Ieremia sau unul din profeți.”


\par }{\PP \VS{15}El le-a zis: “{\WJ{Dar voi cine ziceți că sunt Eu?”
}}


\par }{\PP \VS{16}Simon Petru a răspuns: “Tu ești Hristosul, Fiul lui Dumnezeu cel viu.”


\par }{\PP \VS{17}Isus i-a răspuns:
{\WJ{“Binecuvântat ești tu, Simon Bar Iona, căci nu carnea și sângele ți-au descoperit acest lucru, ci Tatăl Meu care este în ceruri.
}}

\VS{18}{\WJ{Îți spun și Eu îți spun că tu ești Petru și pe această piatră voi zidi Adunarea Mea și porțile Hadesului nu o vor birui.
}}

\VS{19}{\WJ{Îți voi da cheile Împărăției Cerurilor și tot ce vei lega pe pământ va fi fost legat în ceruri; și tot ce vei elibera pe pământ va fi fost eliberat în ceruri.”
}}

\VS{20}Apoi le-a poruncit discipolilor să nu spună nimănui că El este Isus Hristosul.


\par }{\PP \VS{21}Din clipa aceea, Isus a început să le arate ucenicilor Săi că trebuie să meargă la Ierusalim, să sufere multe de la bătrâni, de la preoții cei mai de seamă și de la cărturari, să fie omorât și a treia zi să învie.


\par }{\PP \VS{22}Petru, luându-l deoparte, a început să-L certe, zicând: “Departe de Tine, Doamne! Nu ți se va face niciodată așa ceva”.


\par }{\PP \VS{23}Dar El, întorcându-se, a zis lui Petru: “{\WJ{Satana, treci înapoia mea! Ești o piatră de poticnire pentru mine, pentru că nu te gândești la lucrurile lui Dumnezeu, ci la cele ale oamenilor.”
}}


\par }{\PP \VS{24}Atunci Isus a zis ucenicilor Săi: “{\WJ{Dacă voiește cineva să vină după Mine, să se lepede de sine, să-și ia crucea și să Mă urmeze.
}}

\VS{25}{\WJ{Căci oricine vrea să-și salveze viața o va pierde, iar oricine își va pierde viața pentru Mine o va găsi.
}}

\VS{26}{\WJ{Căci ce-i va folosi cuiva dacă va câștiga toată lumea și-și va pierde viața? Sau ce va da un om în schimbul vieții sale?
}}

\VS{27}{\WJ{Căci Fiul Omului va veni în slava Tatălui Său cu îngerii Săi și atunci va da fiecăruia după faptele sale.
}}

\VS{28}{\WJ{Adevărat vă spun că sunt unii care stau aici și care nu vor gusta nicidecum moartea până când nu vor vedea pe Fiul Omului venind în Împărăția Sa.”
}}



\par }\Chap{17}{\PP \VerseOne{1}După șase zile, Isus a luat cu El pe Petru, pe Iacov și pe Ioan, fratele său, și i-a dus singuri pe un munte înalt.

\VS{2}El s-a schimbat înaintea lor. Fața Lui a strălucit ca soarele, iar hainele Lui au devenit albe ca lumina.

\VS{3}Iată că Moise și Ilie li s-au arătat vorbind cu el.


\par }{\PP \VS{4}Petru a luat cuvântul și a zis lui Isus: “Doamne, bine este să fim aici. Dacă vrei, să facem aici trei corturi: unul pentru tine, unul pentru Moise și unul pentru Ilie.”


\par }{\PP \VS{5}Pe când vorbea El încă, iată că un nor strălucitor i-a acoperit cu umbra sa. Și iată că din nor a ieșit un glas care zicea: “Acesta este Fiul Meu preaiubit, în care Îmi găsesc plăcerea. Ascultați-L!”.


\par }{\PP \VS{6}Ucenicii, când au auzit, au căzut cu fața la pământ și s-au înspăimântat foarte tare.

\VS{7}Isus a venit, i-a atins și le-a zis:
{\WJ{“Ridicați-vă și nu vă temeți.”
}}

\VS{8}Ridicându-și ochii, nu au văzut pe nimeni, decât pe Isus singur.


\par }{\PP \VS{9}Pe când se coborau de pe munte, Isus le-a poruncit:
{\WJ{“Să nu spuneți nimănui ce ați văzut, până ce Fiul Omului nu va învia din morți.”
}}


\par }{\PP \VS{10}Ucenicii Lui L-au întrebat și I-au zis: “Atunci de ce zic cărturarii că Ilie trebuie să vină mai întâi?”


\par }{\PP \VS{11}Isus le-a răspuns: “{\WJ{Ilie vine mai întâi și va restabili toate lucrurile;
}}

\VS{12}{\WJ{dar Eu vă spun că Ilie a venit deja, și ei nu L-au recunoscut, ci i-au făcut ce au vrut. Tot așa și Fiul Omului va suferi din partea lor”.
}}

\VS{13}Atunci discipolii au înțeles că le vorbea despre Ioan Botezătorul.


\par }{\PP \VS{14}Când au ajuns la mulțime, a venit la El un om, care a îngenuncheat înaintea Lui și I-a zis:

\VS{15}“Doamne, ai milă de fiul meu, care este epileptic și suferă cumplit, căci de multe ori cade în foc și de multe ori în apă.

\VS{16}L-am adus la discipolii Tăi și nu au putut să-l vindece.”


\par }{\PP \VS{17}Isus a răspuns:
{\WJ{“Neam necredincios și pervers! Până când voi mai fi cu voi? Până când vă voi răbda? Aduceți-l aici la Mine!”
}}

\VS{18}Isus a mustrat demonul și acesta a ieșit din el, iar băiatul a fost vindecat din ceasul acela.


\par }{\PP \VS{19}Atunci ucenicii au venit la Isus în particular și au zis: “De ce n-am putut să-l scoatem afară?”


\par }{\PP \VS{20}El le-a zis: “Din pricina
{\WJ{necredinței voastre. Căci vă spun cu adevărat că, dacă veți avea credință cât un grăunte de muștar, veți spune acestui munte: “Mută-te de aici încolo”, și se va muta; și nimic nu vă va fi imposibil.
}}

\VS{21}{\WJ{Dar acest fel nu iese decât prin rugăciune și prin post.”
}}


\par }{\PP \VS{22}Pe când stăteau în Galileea, Isus le-a zis: “{\WJ{Fiul Omului va fi dat în mâinile oamenilor
}}

\VS{23}și
{\WJ{ei Îl vor ucide, iar a treia zi va învia.”
}}

\par }{\PP Le-a părut nespus de rău.


\par }{\PP \VS{24}După ce au ajuns în Capernaum, cei care strângeau monedele de didracma au venit la Petru și i-au zis: “Nu cumva învățătorul tău plătește didracma?”

\VS{25}El a răspuns: “Ba da.”

\par }{\PP Când a intrat în casă, Isus l-a anticipat și i-a zis: “{\WJ{Ce crezi, Simone? De la cine primesc regii pământului taxă sau tribut? De la copiii lor sau de la străini?”.
}}


\par }{\PP \VS{26}Petru i-a zis: “De la străini.”

\par }{\PP Isus i-a spus: “{\WJ{De aceea copiii sunt scutiți.
}}

\VS{27}{\WJ{Dar, ca să nu-i facem să se poticnească, du-te la mare, aruncă un cârlig și ia primul pește care iese la suprafață. Când îi vei deschide gura, vei găsi o monedă de un stater. Ia-o și dă-le-o pentru mine și pentru tine.”
}}



\par }\Chap{18}{\PP \VerseOne{1}În ceasul acela, ucenicii au venit la Isus și au zis: “Cine este cel mai mare în Împărăția cerurilor?”


\par }{\PP \VS{2}Isus a chemat la Sine un copilaș, l-a așezat în mijlocul lor

\VS{3}și le-a zis: “{\WJ{Adevărat vă spun că, dacă nu vă veți întoarce și nu veți fi ca niște copilași, nu veți intra nicidecum în Împărăția cerurilor.
}}

\VS{4}{\WJ{Așadar, oricine se smerește ca acest copilaș este cel mai mare în Împărăția Cerurilor.}}

\VS{5}{\WJ{Oricine primește un astfel de copilaș în numele meu mă primește pe mine,
}}

\VS{6}{\WJ{dar oricine face să se poticnească unul dintre acești copilași care cred în mine, ar fi mai bine pentru el dacă i s-ar atârna de gât o piatră de moară uriașă și s-ar scufunda în adâncul mării.
}}


\par }{\PP \VS{7}{\WJ{“Vai de lume, din pricina prilejurilor de poticnire! Căci trebuie să vină ocaziile, dar vai de persoana prin care vine ocazia!
}}

\VS{8}{\WJ{Dacă mâna sau piciorul tău te face să te poticnești, taie-le și leapădă-le de la tine. Este mai bine pentru tine să intri în viață mutilat sau schilodit, decât să ai două mâini sau două picioare pentru a fi aruncat în focul veșnic.
}}

\VS{9}{\WJ{Dacă ochiul tău te face să te poticnești, smulge-l și aruncă-l de la tine. Este mai bine pentru tine să intri în viață cu un singur ochi, decât să ai doi ochi pentru a fi aruncat în gheena de foc.
}}

\VS{10}{\WJ{Ai grijă să nu disprețuiești pe unul dintre acești micuți, căci vă spun că în ceruri îngerii lor văd întotdeauna fața Tatălui meu care este în ceruri.
}}

\VS{11}{\WJ{Căci Fiul Omului a venit să salveze ceea ce era pierdut.
}}


\par }{\PP \VS{12}{\WJ{“Ce părere ai? Dacă un om are o sută de oi și una dintre ele se rătăcește, nu le lasă el pe cele nouăzeci și nouă, nu se duce la munte și nu o caută pe cea rătăcită?
}}

\VS{13}{\WJ{Dacă o găsește, cu siguranță vă spun că se bucură pentru ea mai mult decât pentru cele nouăzeci și nouă care nu s-au rătăcit.
}}

\VS{14}{\WJ{Tot așa, nici Tatăl vostru, care este în ceruri, nu vrea să piară unul dintre acești prunci.
}}


\par }{\PP \VS{15}{\WJ{Dacă fratele tău păcătuiește împotriva ta, du-te și arată-i greșeala lui, numai între tine și el. Dacă te ascultă, ai recâștigat pe fratele tău.
}}

\VS{16}{\WJ{Dar, dacă nu ascultă, ia cu tine încă unul sau doi, pentru ca, la gura a doi sau trei martori, orice cuvânt să fie confirmat.
}}

\VS{17}{\WJ{Dacă refuză să-i asculte, spuneți acest lucru adunării. Dacă refuză să asculte și adunarea, să fie pentru voi ca un neam sau ca un vameș.
}}

\VS{18}{\WJ{Adevărat vă spun că tot ce veți lega pe pământ va fi fost legat în cer și tot ce veți dezlega pe pământ va fi fost dezlegat în cer.
}}

\VS{19}{\WJ{Și iarăși, cu adevărat vă spun că, dacă doi dintre voi se vor învoi pe pământ cu privire la orice lucru pe care îl vor cere, li se va face de către Tatăl Meu care este în ceruri.
}}

\VS{20}{\WJ{Căci unde sunt doi sau trei adunați în numele Meu, acolo sunt și Eu în mijlocul lor.”
}}


\par }{\PP \VS{21}Atunci Petru a venit și I-a zis: “Doamne, de câte ori va păcătui fratele meu împotriva mea, și eu îl voi ierta? Până de șapte ori?”


\par }{\PP \VS{22}Isus i-a zis: “{\WJ{Nu-ți spun până la șapte ori, ci până la șaptezeci de ori câte șapte.
}}

\VS{23}De
{\WJ{aceea, Împărăția cerurilor se aseamănă cu un rege care voia să regleze conturile cu slujitorii săi.
}}

\VS{24}{\WJ{Când a început să regleze, i s-a adus unul care îi datora zece mii de talanți.
}}

\VS{25}{\WJ{Dar, pentru că nu putea plăti, stăpânul său a poruncit să fie vândut, împreună cu soția, copiii și tot ce avea, și să se facă plata.
}}

\VS{26}{\WJ{Slujitorul a căzut, așadar, și a îngenuncheat înaintea lui, zicând: “Doamne, ai răbdare cu mine și îți voi plăti totul!”.
}}

\VS{27}{\WJ{Stăpânul acelui rob, cuprins de compasiune, i-a dat drumul și i-a iertat datoria.
}}


\par }{\PP \VS{28}{\WJ{Dar robul acela a ieșit afară și a găsit pe unul din tovarășii săi de slujbă care îi datora o sută de denari; l-a apucat și l-a luat de gât, zicând: “Plătește-mi ce ai de plătit!”.
}}


\par }{\PP \VS{29}{\WJ{Și a căzut sluga lui la picioarele lui și l-a rugat, zicând: “Ai răbdare cu mine și-ți voi răsplăti!”.
}}

\VS{30}Dar
{\WJ{el n-a vrut, ci s-a dus și l-a aruncat în temniță până ce va plăti ce i se cuvine.
}}

\VS{31}Când au
{\WJ{văzut tovarășii lui de slujbă ce se făcuse, s-au mâhnit foarte tare și au venit și au povestit stăpânului lor tot ce se făcuse.
}}

\VS{32}{\WJ{Atunci stăpânul său l-a chemat înăuntru și i-a zis: “Slugă ticăloasă! Ți-am iertat toată datoria aceea pentru că m-ai implorat.
}}

\VS{33}N-ar fi trebuit să ai
{\WJ{și tu milă de tovarășul tău de slujbă, așa cum am avut eu milă de tine?”}}

\VS{34}{\WJ{Stăpânul său s-a mâniat și l-a dat pe el torționarilor până când va plăti tot ce i se cuvenea.
}}

\VS{35}{\WJ{Așava face și Tatăl meu ceresc cu voi, dacă nu-l veți ierta fiecare din inimă pe fratele vostru pentru greșelile lui.”
}}



\par }\Chap{19}{\PP \VerseOne{1}După ce a isprăvit Isus aceste cuvinte, a plecat din Galileea și a ajuns în ținutul Iudeii, dincolo de Iordan.

\VS{2}L-au urmat mulțimi mari de oameni și El i-a vindecat acolo.


\par }{\PP \VS{3}Fariseii au venit la El și L-au pus la încercare, zicând: “Este îngăduit oare unui bărbat să divorțeze de nevastă-sa pentru orice motiv?”


\par }{\PP \VS{4}El a răspuns: “{\WJ{Nu ați citit că Cel ce i-a făcut de la început i-a făcut bărbat și femeie,
}}

\VS{5}{\WJ{și a zis: “De aceea va lăsa omul pe tatăl său și pe mama sa și se va uni cu femeia sa, și cei doi vor fi un singur trup?”
}}

\VS{6}{\WJ{Astfel că nu mai sunt doi, ci un singur trup. Așadar, ceea ce Dumnezeu a unit, să nu lase omul să despartă.”
}}


\par }{\PP \VS{7}Și l-au întrebat: “Atunci de ce ne-a poruncit Moise să-i dăm certificatul de divorț și să divorțăm de ea?”


\par }{\PP \VS{8}Și le-a zis: “{\WJ{Moise, din pricina împietririi inimilor voastre, v-a îngăduit să divorțați de nevestele voastre, dar de la început nu a fost așa.
}}

\VS{9}{\WJ{Eu vă spun că oricine divorțează de soția sa, cu excepția imoralității sexuale, și se căsătorește cu alta, comite adulter; iar cel care se căsătorește cu ea când este divorțată comite adulter.”
}}


\par }{\PP \VS{10}Ucenicii Lui I-au zis: “Dacă așa se întâmplă cu bărbatul și cu nevasta lui, nu este bine să te căsătorești.”


\par }{\PP \VS{11}Dar El le-a zis: “{\WJ{Nu toți oamenii pot primi cuvântul acesta, ci numai aceia cărora le este dat.
}}

\VS{12}{\WJ{Căci sunt eunuci care s-au născut așa din pântecele mamei lor, sunt eunuci care au fost făcuți eunuci de către oameni și sunt eunuci care s-au făcut eunuci de dragul Împărăției Cerurilor. Cine este în stare să o primească, să o primească”.
}}


\par }{\PP \VS{13}Atunci I-au fost aduși niște copilași, ca să pună mâinile peste ei și să se roage; dar ucenicii îi mustrau.

\VS{14}Dar Isus a zis: “{\WJ{Lăsați-i pe copilași și nu le interziceți să vină la Mine, căci Împărăția Cerurilor este a unora ca aceștia.”
}}

\VS{15}Și-a pus mâinile peste ei și a plecat de acolo.


\par }{\PP \VS{16}Și iată că a venit cineva la El și a zis: “Învățătorule, ce să fac ca să am viața veșnică?”


\par }{\PP \VS{17}El i-a zis: “{\WJ{De ce Mă numești bun? Nimeni nu este bun decât unul singur, adică Dumnezeu. Dar dacă vrei să intri în viață, păzește poruncile”.
}}


\par }{\PP \VS{18}El i-a zis: “Care dintre ele?”

\par }{\PP Isus a spus: “{\WJ{Să nu ucizi. Să nu comiți adulter”. 'Să nu furi'. 'Să nu depui mărturie mincinoasă'.
}}

\VS{19}{\WJ{'Cinstește pe tatăl tău și pe mama ta'. Și: 'Să iubești pe aproapele tău ca pe tine însuți'”.
}}


\par }{\PP \VS{20}Tânărul i-a zis: “Toate acestea le știu din tinerețea mea. Ce-mi mai lipsește?”


\par }{\PP \VS{21}Isus i-a zis: “{\WJ{Dacă vrei să fii desăvârșit, du-te, vinde ce ai, dă săracilor și vei avea o comoară în ceruri; apoi vino și urmează-Mi.”
}}

\VS{22}Dar tânărul, când a auzit aceasta, a plecat trist, pentru că era unul care avea averi mari.


\par }{\PP \VS{23}Isus a zis ucenicilor Săi: “{\WJ{Adevărat vă spun că un om bogat va intra cu greu în Împărăția cerurilor.
}}

\VS{24}{\WJ{Vă spun din nou: este mai ușor să treacă o cămilă prin urechea unui ac decât să intre un bogat în Împărăția lui Dumnezeu.”
}}


\par }{\PP \VS{25}Ucenicii, auzind aceasta, au rămas foarte mirați, și au zis: “Cine poate fi mântuit?”


\par }{\PP \VS{26}Și, uitându-se la ei, Isus a zis: “La
{\WJ{oameni este cu neputință, dar la Dumnezeu toate sunt cu putință.”
}}


\par }{\PP \VS{27}Atunci Petru a răspuns: “Iată, am lăsat totul și am venit după Tine. Ce vom avea atunci?”


\par }{\PP \VS{28}Isus le-a zis: “{\WJ{Adevărat vă spun că voi, cei ce M-ați urmat, în vremea când Fiul Omului va ședea pe scaunul de domnie al slavei Sale, veți ședea și voi pe douăsprezece scaune de domnie, ca să judecați cele douăsprezece seminții ale lui Israel.
}}

\VS{29}{\WJ{Oricine a lăsat case, sau frați, sau surori, sau tată, sau mamă, sau soție, sau copii, sau pământuri, de dragul numelui Meu, va primi de o sută de ori și va moșteni viața veșnică.
}}

\VS{30}{\WJ{Dar mulți vor fi cei din urmă care sunt primii și cei dintâi care sunt ultimii.
}}



\par }\Chap{20}{\PP \VerseOne{1}{\WJ{“Căci Împărăția cerurilor se aseamănă cu un om, stăpânul unei gospodării, care a ieșit dis-de-dimineață să angajeze lucrători pentru via sa.
}}

\VS{2}După ce s-a
{\WJ{înțeles cu muncitorii pentru un denar pe zi, i-a trimis în via sa.
}}

\VS{3}{\WJ{A ieșit pe la ora trei și a văzut pe alții care stăteau degeaba în piață.
}}

\VS{4}{\WJ{Și le-a zis: “Mergeți și voi în vie și vă voi da ce este drept. Și ei au plecat.
}}

\VS{5}{\WJ{A ieșit iarăși pe la ceasul al șaselea și al nouălea și a făcut la fel.
}}

\VS{6}{\WJ{Pe la ceasul al unsprezecelea a ieșit și a găsit pe alții care stăteau degeaba. El i-a întrebat: “De ce stați aici toată ziua degeaba?}}”.


\par }{\PP \VS{7}{\WJ{Ei I-au răspuns: “Pentru că nimeni nu ne-a angajat.
}}

\par }{\PP {\WJ{“El le-a zis: “Mergeți și voi în vie și veți primi tot ce este drept”.
}}


\par }{\PP \VS{8}{\WJ{“Când s-a făcut seară, stăpânul viei a zis administratorului său: “Cheamă lucrătorii și dă-le plata, de la cel din urmă până la cel dintâi”.
}}

\VS{9}{\WJ{“Când au venit cei care fuseseră angajați pe la ora unsprezece, au primit fiecare câte un denar.
}}

\VS{10}{\WJ{Când au venit cei dintâi, au presupus că vor primi mai mult; și, de asemenea, și ei au primit fiecare câte un denar.
}}

\VS{11}{\WJ{Când l-au primit, au cârtit împotriva stăpânului casei,
}}

\VS{12}{\WJ{zicând: “Aceștia din urmă au cheltuit un singur ceas, și tu i-ai făcut egali cu noi, care am dus povara zilei și căldura arzătoare!}}”.


\par }{\PP \VS{13}{\WJ{Iar el a răspuns unuia dintre ei: “Prietene, nu-ți fac nici un rău. Nu v-ați înțeles cu mine pentru un denar?
}}

\VS{14}{\WJ{Ia ceea ce este al tău și du-te. Dorința mea este să dau și acestui din urmă la fel de mult ca și ție.
}}

\VS{15}Nu-mi este
{\WJ{oare permis să fac ce vreau cu ceea ce am? Sau ochiul tău este rău, pentru că eu sunt bun?”
}}

\VS{16}{\WJ{Astfel, cel din urmă va fi cel dintâi, iar cel dintâi cel din urmă. Căci mulți sunt chemați, dar puțini sunt aleși.”
}}


\par }{\PP \VS{17}Pe când se suia Isus la Ierusalim, i-a luat deoparte pe cei doisprezece ucenici și, pe drum, le-a zis:

\VS{18}“{\WJ{Iată, ne suim la Ierusalim, și Fiul Omului va fi dat preoților de seamă și cărturarilor, care Îl vor condamna la moarte,
}}

\VS{19}{\WJ{și-L vor da neamurilor ca să-L batjocorească, să-L biciuiască și să-L răstignească, iar a treia zi va învia.”
}}


\par }{\PP \VS{20}Atunci mama fiilor lui Zebedei a venit la El cu fiii ei, îngenunchind și cerându-I un lucru.

\VS{21}El a întrebat-o:
{\WJ{“Ce vrei?”
}}

\par }{\PP Ea i-a zis: “Poruncește ca acești doi fii ai mei să șadă, unul la dreapta Ta și altul la stânga Ta, în Împărăția Ta”.


\par }{\PP \VS{22}Dar Isus i-a răspuns: “{\WJ{Nu știi ce ceri. Ești în stare să bei paharul pe care voi bea eu și să te botezi cu botezul cu care mă botez eu?”
}}

\par }{\PP Ei i-au spus: “Noi putem”.


\par }{\PP \VS{23}El le-a zis: “{\WJ{Voi veți bea paharul Meu și veți fi botezați cu botezul cu care sunt botezat Eu; dar șederea la dreapta și la stânga Mea nu este a Mea, ci a cui a fost pregătită de Tatăl Meu.”
}}


\par }{\PP \VS{24}Când au auzit cei zece, s-au mâniat pe cei doi frați.


\par }{\PP \VS{25}Dar Isus i-a chemat și le-a zis: “{\WJ{Știți că domnitorii neamurilor le stăpânesc și că cei mari ale lor le stăpânesc.}}

\VS{26}{\WJ{Nu va fi așa între voi; ci oricine vrea să devină mare între voi va fi robul vostru.
}}

\VS{27}{\WJ{Oricine vrea să fie cel dintâi între voi va fi robul vostru,
}}

\VS{28}{\WJ{așa cum Fiul Omului nu a venit ca să fie slujit, ci ca să slujească și să-și dea viața ca răscumpărare pentru mulți.”
}}


\par }{\PP \VS{29}Când au ieșit din Ierihon, o mare mulțime a mers după El.

\VS{30}Și iată că doi orbi care stăteau pe drum, auzind că Isus trece pe acolo, au strigat: “Doamne, ai milă de noi, Fiul lui David!”

\VS{31}Mulțimea i-a mustrat, spunându-le să tacă, dar ei au strigat și mai tare: “Doamne, ai milă de noi, Fiul lui David!”


\par }{\PP \VS{32}Isus, oprindu-se, i-a chemat și i-a întrebat:
{\WJ{“Ce vreți să fac pentru voi?”
}}


\par }{\PP \VS{33}Și I-au zis: “Doamne, ca să ni se deschidă ochii.”


\par }{\PP \VS{34}Isus, cuprins de compasiune, s-a atins de ochii lor; și îndată ochii lor și-au recăpătat vederea, și au mers după El.



\par }\Chap{21}{\PP \VerseOne{1}Când s-au apropiat de Ierusalim și au ajuns la Betfaghe, pe Muntele Măslinilor, Isus a trimis doi ucenici,

\VS{2}și le-a zis: “{\WJ{Mergeți în satul care este în fața voastră și îndată veți găsi o măgăriță legată și un mânz cu ea. Dezlegați-i și aduceți-i la mine.
}}

\VS{3}{\WJ{Dacă vă va spune cineva ceva, să spuneți: “Domnul are nevoie de ei”, și îndată îi va trimite.”
}}


\par }{\PP \VS{4}Toate acestea s-au făcut pentru ca să se împlinească ceea ce fusese spus prin proorocul care zice,


\par }{\Q \VS{5}“Spune fiicei Sionului,

\par }{\QB Iată, Regele vostru vine la voi,

\par }{\QB umil, și călare pe un măgar,

\par }{\QB pe un mânz, mânzul unei măgărițe.”


\par }{\PP \VS{6}Ucenicii s-au dus și au făcut cum le poruncise Isus;

\VS{7}au adus măgărița și mânzul, au pus hainele lor pe ei și Isus a șezut pe ei.

\VS{8}O mulțime foarte mare și-a întins hainele pe drum. Alții au tăiat ramuri din copaci și le-au întins pe drum.

\VS{9}Mulțimea care mergea în fața lui și cei care îl urmau strigau întruna: “Osana Fiului lui David! Binecuvântat este cel care vine în numele Domnului! Osana în cele mai înalte!”


\par }{\PP \VS{10}După ce a intrat în Ierusalim, toată cetatea s-a tulburat și zicea: “Cine este acesta?”


\par }{\PP \VS{11}Și mulțimile ziceau: “Acesta este proorocul Isus, din Nazaret din Galileea.”


\par }{\PP \VS{12}Isus a intrat în Templul lui Dumnezeu și a izgonit pe toți cei ce vindeau și cumpărau în Templu, a răsturnat mesele schimbătorilor de bani și scaunele celor ce vindeau porumbei.

\VS{13}El le-a zis:
{\WJ{“Este scris: “Casa Mea se va numi casă de rugăciune”, dar voi ați făcut din ea o peșteră de tâlhari!”
}}


\par }{\PP \VS{14}Șchiopii și orbii veneau la El în templu, și El îi vindeca.

\VS{15}Dar preoții cei mai de seamă și cărturarii, văzând minunile pe care le făcea și pe copiii care strigau în templu și ziceau: “Osana Fiului lui David!”, s-au indignat

\VS{16}și i-au zis: “Auzi ce spun aceștia?”

\par }{\PP Isus le-a zis:
{\WJ{“Da, dar nu ați citit niciodată: “Din gura copiilor și a celor ce alăptează, ai desăvârșit lauda?””.
}}


\par }{\PP \VS{17}Și lăsându-i, a ieșit din cetate și s-a dus în Betania și a tăbărât acolo.


\par }{\PP \VS{18}Dimineața, pe când se întorcea în cetate, îi era foame.

\VS{19}Văzând un smochin lângă drum, s-a apropiat de el și n-a găsit pe el decât frunze. I-a spus:
{\WJ{“Să nu mai fie niciun fruct de la tine în veci!”
}}

\par }{\PP Imediat smochinul s-a uscat.


\par }{\PP \VS{20}Ucenicii, văzând aceasta, s-au mirat și au zis: “Cum de s-a uscat smochinul îndată?”


\par }{\PP \VS{21}Isus le-a răspuns: “{\WJ{Adevărat vă spun că, dacă aveți credință și nu vă îndoiți, nu numai că veți face ce s-a făcut cu smochinul, ci chiar dacă ați spune muntelui acesta: “Ridică-te și aruncă-te în mare”, se va face.
}}

\VS{22}{\WJ{Toate lucrurile, orice veți cere în rugăciune, crezând, le veți primi.”
}}


\par }{\PP \VS{23}După ce a intrat în Templu, preoții cei mai de seamă și bătrânii norodului au venit la El, pe când învăța El, și I-au zis: “Cu ce putere faci Tu aceste lucruri? Cine ți-a dat această autoritate?”


\par }{\PP \VS{24}Isus le-a răspuns: “{\WJ{Și eu vă voi pune o întrebare; dacă îmi veți răspunde, vă voi spune și eu cu ce putere fac aceste lucruri.
}}

\VS{25}{\WJ{Botezul lui Ioan, de unde a fost? Din cer sau de la oameni?”
}}

\par }{\PP Ei se gândeau între ei, zicând: “Dacă spunem: “Din cer”, ne va întreba: “De ce nu l-ați crezut?”

\VS{26}Dar dacă spunem: “De la oameni”, ne temem de mulțime, căci toți îl consideră pe Ioan ca profet”.

\VS{27}Ei au răspuns lui Isus și au zis: “Nu știm.”

\par }{\PP El le-a mai spus: “{\WJ{Nu vă voi spune nici cu ce putere fac aceste lucruri.
}}

\VS{28}{\WJ{Dar voi ce credeți? Un om avea doi fii; a venit la primul și i-a zis: “Fiule, du-te astăzi să lucrezi în via mea!”.
}}

\VS{29}{\WJ{El a răspuns: “Nu vreau!” Dar, după aceea, s-a răzgândit și s-a dus.
}}

\VS{30}{\WJ{A venit la al doilea și i-a spus același lucru. El a răspuns: “Mă duc, domnule”, dar nu s-a dus.
}}

\VS{31}{\WJ{Care dintre cei doi a făcut voia tatălui său?”
}}

\par }{\PP I-au spus: “Cel dintâi”.

\par }{\PP Isus le-a zis: “{\WJ{Adevărat vă spun că vameșii și prostituatele intră înaintea voastră în Împărăția lui Dumnezeu.
}}

\VS{32}{\WJ{Căci Ioan a venit la voi pe calea dreptății și voi nu l-ați crezut, dar vameșii și prostituatele l-au crezut. Când l-ați văzut, nici măcar nu v-ați pocăit după aceea, ca să-l credeți.
}}


\par }{\PP \VS{33}{\WJ{“Ascultați o altă pildă. Era un om care era stăpânul unei gospodării și care a plantat o vie, a pus un gard în jurul ei, a săpat în ea o presă de vin, a construit un turn, a dat-o în arendă unor agricultori și a plecat în altă țară.
}}

\VS{34}{\WJ{Când s-a apropiat sezonul de recoltare a roadei, el și-a trimis slujitorii la fermieri ca să-i primească roadele.
}}

\VS{35}{\WJ{Fermierii i-au luat slujitorii, au bătut pe unul, au ucis pe altul și au ucis cu pietre pe altul.
}}

\VS{36}{\WJ{A trimis din nou alți slujitori, mai mulți decât primii, și s-au purtat cu ei la fel.
}}

\VS{37}{\WJ{După aceea însă, a trimis la ei pe fiul său, spunând: “Îl vor respecta pe fiul meu”.
}}

\VS{38}{\WJ{Dar țăranii, când au văzut pe fiu, au zis între ei: “Acesta este moștenitorul”. Haideți să-l omorâm și să-i luăm moștenirea'.
}}

\VS{39}{\WJ{Și l-au luat și l-au aruncat afară din vie, apoi l-au omorât.
}}

\VS{40}{\WJ{Când va veni, așadar, stăpânul viei, ce va face cu acei țărani?”
}}


\par }{\PP \VS{41}Ei i-au spus: “Îi va nimici pe acești oameni nenorociți și va da via în arendă altor agricultori, care îi vor da roadele la vremea lor.”


\par }{\PP \VS{42}Isus le-a zis: “{\WJ{N-ați citit niciodată în Scriptură?
}}

\par }{\Q {\WJ{“Piatra pe care au respins-o zidarii
}}

\par }{\QB {\WJ{a fost numit șef de colț.
}}

\par }{\Q {\WJ{Aceasta a fost de la Domnul.
}}

\par }{\QB {\WJ{Este minunat în ochii noștri”?
}}


\par }{\PP \VS{43}De
{\WJ{aceea vă spun că Împărăția lui Dumnezeu va fi luată de la voi și va fi dată unui neam care își va face roadele.
}}

\VS{44}Cine va cădea
{\WJ{pe piatra aceasta va fi sfărâmat, dar pe oricine va cădea, îl va împrăștia ca praful.”
}}


\par }{\PP \VS{45}Preoții cei mai de seamă și fariseii, auzind pildele Lui, au înțeles că vorbea despre ei.

\VS{46}Când au căutat să îl prindă, s-au temut de mulțime, pentru că îl considerau profet.



\par }\Chap{22}{\PP \VerseOne{1}Isus, răspunzând, le-a vorbit iarăși în pilde, zicând:

\VS{2}{\WJ{“Împărăția cerurilor se aseamănă cu un împărat care a făcut un ospăț de nuntă pentru fiul său
}}

\VS{3}și a
{\WJ{trimis pe slujitorii săi să cheme pe cei invitați la nuntă, dar aceștia n-au vrut să vină.
}}

\VS{4}{\WJ{A trimis din nou alți slujitori, spunând: “Spuneți-le celor care sunt invitați: “Iată, am pregătit cina mea. Vitele și îngrășămintele mele sunt ucise și toate lucrurile sunt gata. Veniți la ospățul de nuntă!””'”.
}}

\VS{5}{\WJ{Dar ei au luat totul în derâdere și au plecat, unul la ferma lui, altul la marfa lui;
}}

\VS{6}{\WJ{iar ceilalți i-au înșfăcat slujitorii, i-au tratat cu rușine și i-au ucis.
}}

\VS{7}{\WJ{Când a auzit asta, regele s-a înfuriat și și-a trimis armatele, i-a nimicit pe acei ucigași și le-a ars cetatea.
}}


\par }{\PP \VS{8}{\WJ{Și a zis robilor săi: “Nunta este gata, dar cei ce au fost invitați nu au fost vrednici.
}}

\VS{9}{\WJ{Duceți-vă, așadar, la intersecțiile drumurilor și, pe câți veți găsi, invitați-i la nuntă'.
}}

\VS{10}{\WJ{Slujitorii aceia au ieșit pe drumuri și au adunat pe câți au găsit, atât pe cei răi, cât și pe cei buni. Nunta s-a umplut de invitați.
}}


\par }{\PP \VS{11}{\WJ{Dar când a intrat regele să vadă oaspeții, a văzut acolo un om care nu avea haine de nuntă,
}}

\VS{12}{\WJ{și l-a întrebat: “Prietene, cum ai intrat aici fără haine de nuntă?” El a rămas fără cuvinte.
}}

\VS{13}{\WJ{Atunci regele a zis slujitorilor: “Legați-l de mâini și de picioare, luați-l și aruncați-l în întunericul de afară. Acolo va fi plânsul și scrâșnitul dinților”.
}}

\VS{14}{\WJ{Căci mulți sunt chemați, dar puțini sunt aleși.”
}}


\par }{\PP \VS{15}Atunci fariseii s-au dus și s-au sfătuit cum să-l prindă în discuția lui.

\VS{16}Și au trimis la el pe ucenicii lor și pe Irodieni, zicând: “Învățătorule, știm că ești cinstit și că înveți calea lui Dumnezeu cu adevărat, indiferent pe cine înveți, pentru că nu ești părtinitor față de nimeni.

\VS{17}Spuneți-ne, așadar, ce credeți? Este sau nu legitim să plătim impozite Cezarului?”


\par }{\PP \VS{18}Dar Isus, înțelegând răutatea lor, a zis: “{\WJ{Pentru ce Mă încercați, fățarnicilor?
}}

\VS{19}{\WJ{Arătați-Mi banii de impozit.”
}}

\par }{\PP I-au adus un denar.


\par }{\PP \VS{20}Și i-a întrebat:
{\WJ{“Al cui este chipul acesta și inscripția aceasta?”
}}


\par }{\PP \VS{21}Ei i-au zis: “Al Cezarului”.

\par }{\PP Apoi le-a zis: “{\WJ{Dați deci Cezarului ce este al Cezarului și lui Dumnezeu ce este al lui Dumnezeu”.
}}


\par }{\PP \VS{22}Când au auzit, s-au mirat, L-au lăsat și au plecat.


\par }{\PP \VS{23}În ziua aceea au venit la El saducheii, care zic că nu există înviere. L-au întrebat

\VS{24}și i-au zis: “Învățătorule, Moise a zis: “Dacă un om moare fără să aibă copii, fratele său să se căsătorească cu femeia lui și să ridice urmași fratelui său.

\VS{25}Or, erau cu noi șapte frați. Cel dintâi s-a căsătorit și a murit și, neavând urmași, și-a lăsat soția fratelui său.

\VS{26}La fel, și al doilea, și al treilea, până la al șaptelea.

\VS{27}După ei toți, femeia a murit.

\VS{28}La înviere, așadar, a cui va fi soția dintre cei șapte? Căci toți au avut-o.”


\par }{\PP \VS{29}Dar Isus le-a răspuns: “{\WJ{Vă înșelați, căci nu cunoașteți Scripturile și nu știți puterea lui Dumnezeu.
}}

\VS{30}{\WJ{Căci la înviere nu se căsătoresc și nici nu sunt dați în căsătorie, ci sunt ca îngerii lui Dumnezeu în ceruri.
}}

\VS{31}{\WJ{Dar, în ce privește învierea morților, nu ați citit ce v-a spus Dumnezeu, când a zis:
}}

\VS{32}“{\WJ{Eu sunt Dumnezeul lui Avraam, Dumnezeul lui Isaac și Dumnezeul lui Iacov”? Dumnezeu nu este Dumnezeul morților, ci al celor vii.”
}}


\par }{\PP \VS{33}Când au auzit mulțimile, au rămas uimite de învățătura lui.


\par }{\PP \VS{34}Dar fariseii, auzind că i-a redus la tăcere pe saduchei, s-au adunat.

\VS{35}Unul dintre ei, un avocat, i-a pus o întrebare, punându-l la încercare.

\VS{36}“Învățătorule, care este cea mai mare poruncă din Lege?”


\par }{\PP \VS{37}Isus i-a zis: “{\WJ{Să iubești pe Domnul Dumnezeul tău din toată inima ta, din tot sufletul tău și din tot cugetul tău.
}}

\VS{38}{\WJ{Aceasta este cea dintâi și cea mai mare poruncă.
}}

\VS{39}{\WJ{A doua, la fel, este aceasta: “Să iubești pe aproapele tău ca pe tine însuți”.
}}

\VS{40}{\WJ{Toată Legea și profeții depind de aceste două porunci.”
}}


\par }{\PP \VS{41}Pe când se adunaseră fariseii, Isus le-a pus o întrebare,

\VS{42}și le-a zis:
{\WJ{“Ce credeți voi despre Hristos? Al cui fiu este El?”
}}

\par }{\PP I-au spus: “De David”.


\par }{\PP \VS{43}El le-a zis: “{\WJ{Cum deci David, în Duhul Sfânt, Îl numește Domn, zicând,
}}


\par }{\Q \VS{44}{\WJ{“Domnul a zis Domnului meu,
}}

\par }{\QB {\WJ{stai la dreapta mea,
}}

\par }{\QB {\WJ{până când voi face din vrăjmașii tăi un scăunel pentru picioarele tale”?
}}


\par }{\PP \VS{45}{\WJ{“Dacă David îl numește Domn, cum este el fiul său?”
}}


\par }{\PP \VS{46}Nimeni n-a putut să-i răspundă nimic și nimeni n-a îndrăznit să-i mai pună vreo întrebare din ziua aceea.



\par }\Chap{23}{\PP \VerseOne{1}Atunci Isus a vorbit mulțimilor și ucenicilor Săi,

\VS{2}și a zis: “{\WJ{Cărturarii și fariseii stau pe scaunul lui Moise.
}}

\VS{3}{\WJ{Așadar, tot ce vă spun ei să observați, observați și faceți, dar nu faceți faptele lor; căci ei spun și nu fac.
}}

\VS{4}{\WJ{Căci ei leagă poveri grele, greu de purtat, și le pun pe umerii oamenilor, dar ei înșiși nu vor să ridice un deget pentru a-i ajuta.
}}

\VS{5}{\WJ{Ci își fac toate faptele lor ca să fie văzute de oameni. Își lărgesc filacteriile și își lărgesc franjurile hainelor lor,
}}

\VS{6}{\WJ{și le place locul de cinste la sărbători, locurile cele mai bune din sinagogi,
}}

\VS{7}{\WJ{salutul în piețe și să fie numiți de oameni “Rabi, Rabi, Rabi”.
}}

\VS{8}{\WJ{Dar voi nu trebuie să fiți numiți “Rabi”, căci unul singur este învățătorul vostru, Hristosul, și voi toți sunteți frați.
}}

\VS{9}{\WJ{Nu numiți pe nimeni de pe pământ Tatăl vostru, căci unul singur este Tatăl vostru, Cel din ceruri.
}}

\VS{10}Și nu vă
{\WJ{numiți nici stăpâni, căci unul singur este stăpânul vostru, Hristosul.
}}

\VS{11}{\WJ{Ci cel care este cel mai mare dintre voi va fi robul vostru.
}}

\VS{12}{\WJ{Cine se înalță va fi smerit, iar cine se smerește va fi înălțat.
}}


\par }{\PP \VS{13}“{\WJ{Vai vouă, cărturari și farisei, fățarnici! Căci voi devorați casele văduvelor și, ca un pretext, faceți rugăciuni lungi. De aceea veți primi o condamnare și mai mare.
}}


\par }{\PP \VS{14}“{\WJ{Dar vai de voi, cărturari și farisei, fățarnici! Pentru că voi închideți Împărăția Cerurilor împotriva oamenilor; pentru că voi înșivă nu intrați înăuntru, și nici pe cei care intră nu-i lăsați să intre.
}}

\VS{15}{\WJ{Vai vouă, cărturari și farisei, ipocriților! Pentru că voi călătoriți pe mare și pe uscat ca să faceți un singur proselit; și când acesta devine unul, îl faceți de două ori mai fiu al Gheenei decât voi înșivă.
}}


\par }{\PP \VS{16}“{\WJ{Vai de voi, călăuze oarbe, care ziceți: “Oricine se jură pe templu, nu este nimic; dar oricine se jură pe aurul templului, este obligat.
}}

\VS{17}{\WJ{Voi, nebuni orbi! Căci ce este mai mare, aurul sau templul care sfințește aurul?
}}

\VS{18}{\WJ{Și: “Oricine se jură pe altar, nu este nimic; dar oricine se jură pe darul care este pe el, este obligat?”
}}

\VS{19}{\WJ{Nebuni orbi! Căci ce este mai mare, darul sau altarul care sfințește darul?
}}

\VS{20}{\WJ{Așadar, cine jură pe altar, jură pe el și pe tot ce este pe el.
}}

\VS{21}Cine se
{\WJ{jură pe templu, se jură pe el și pe cel care a locuit în el.
}}

\VS{22}{\WJ{Cine se jură pe cer, se jură pe tronul lui Dumnezeu și pe cel care stă pe el.
}}


\par }{\PP \VS{23}“{\WJ{Vai de voi, cărturari și farisei, fățarnici! Pentru că voi dați zeciuială la mentă, mărar și chimen și ați lăsat deoparte cele mai grele lucruri ale Legii: dreptatea, mila și credința. Dar trebuia să le faceți pe acestea și să nu le lăsați pe celelalte nefăcute.
}}

\VS{24}{\WJ{Voi, călăuze oarbe, care strecurați un țânțar și înghițiți o cămilă!
}}


\par }{\PP \VS{25}“{\WJ{Vai vouă, cărturari și farisei, fățarnici! Pentru că voi curățați partea exterioară a paharului și a farfuriei, dar înăuntru sunt pline de jaf și de nedreptate.
}}

\VS{26}{\WJ{Fariseu orb, curăță mai întâi interiorul paharului și al platoului, pentru ca și exteriorul lui să devină curat.
}}


\par }{\PP \VS{27}“{\WJ{Vai de voi, cărturari și farisei, fățarnici! Pentru că sunteți ca niște morminte albite, care în afară par frumoase, dar înăuntru sunt pline de oase de morți și de toată necurăția.
}}

\VS{28}{\WJ{Tot așa și voi, în aparență, păreți drepți în fața oamenilor, dar pe dinăuntru sunteți plini de fățărnicie și de nelegiuire.
}}


\par }{\PP \VS{29}“{\WJ{Vai vouă, cărturari și farisei, fățarnici! Căci voi zidiți mormintele profeților și împodobiți mormintele drepților
}}

\VS{30}{\WJ{și spuneți: “Dacă am fi trăit în zilele părinților noștri, nu am fi fost părtași cu ei la sângele profeților”.
}}

\VS{31}{\WJ{De aceea vă mărturisiți vouă înșivă că sunteți copiii celor care au ucis pe profeți.
}}

\VS{32}{\WJ{Umpleți, așadar, măsura părinților voștri.
}}

\VS{33}{\WJ{Voi, șerpi, urmași ai viperelor, cum veți scăpa de judecata Gheenei?
}}

\VS{34}{\WJ{De aceea, iată, trimit la voi profeți, înțelepți și cărturari. Pe unii dintre ei îi veți ucide și îi veți răstigni, iar pe alții îi veți biciui în sinagogile voastre și îi veți persecuta din cetate în cetate,
}}

\VS{35}pentru
{\WJ{ca asupra voastră să cadă tot sângele neprihănit vărsat pe pământ, de la sângele neprihănitului Abel până la sângele lui Zaharia, fiul lui Barachia, pe care l-ați ucis între sanctuar și altar.
}}

\VS{36}Cu
{\WJ{siguranță vă spun că toate aceste lucruri vor veni peste această generație.
}}


\par }{\PP \VS{37}{\WJ{“Ierusalime, Ierusalime, care ucizi pe prooroci și ucizi cu pietre pe cei trimiși la ea! De câte ori aș fi vrut să-ți adun copiii, cum își adună găina puii sub aripi, și tu n-ai vrut!
}}

\VS{38}{\WJ{Iată, casa ta este lăsată pustie pentru tine.
}}

\VS{39}{\WJ{Căcivă spun că nu mă veți mai vedea de acum înainte, până când nu veți zice: “Binecuvântat este cel ce vine în numele Domnului!”.”
}}



\par }\Chap{24}{\PP \VerseOne{1}Isus a ieșit din templu și mergea pe drumul Său. Discipolii săi au venit la el ca să-i arate clădirile templului.

\VS{2}Dar el le-a răspuns: “{\WJ{Voi vedeți toate aceste lucruri, nu-i așa? Cu siguranță vă spun că nu va rămâne aici piatră peste piatră care să nu fie dărâmată.”
}}


\par }{\PP \VS{3}Pe când ședea El pe Muntele Măslinilor, ucenicii au venit la El în particular și I-au zis: “Spune-ne, când se vor întâmpla aceste lucruri? Care este semnul venirii Tale și al sfârșitului veacului?”


\par }{\PP \VS{4}Isus le-a răspuns: “Luați
{\WJ{seama ca nimeni să nu vă rătăcească.
}}

\VS{5}{\WJ{Căci mulți vor veni în numele Meu, zicând: “Eu sunt Hristosul”, și vor duce pe mulți în rătăcire.
}}

\VS{6}{\WJ{Veți auzi de războaie și de zvonuri de războaie. Aveți grijă să nu vă tulburați, căci toate acestea trebuie să se întâmple, dar sfârșitul nu este încă.
}}

\VS{7}{\WJ{Căci se va ridica națiune împotriva națiunii și regat împotriva regatului, și vor fi foamete, ciumă și cutremure în diferite locuri.
}}

\VS{8}{\WJ{Dar toate aceste lucruri sunt începutul durerilor nașterii.
}}


\par }{\PP \VS{9}{\WJ{“Atunci vă vor da pe voi la asuprire și vă vor ucide. Veți fi urâți de toate neamurile din cauza numelui Meu.
}}

\VS{10}{\WJ{Atunci mulți se vor poticni, se vor preda unii pe alții și se vor urî unii pe alții.
}}

\VS{11}Se
{\WJ{vor ridica mulți profeți falși și vor duce pe mulți în rătăcire.
}}

\VS{12}{\WJ{Pentru că se va înmulți nelegiuirea, dragostea multora se va răci.
}}

\VS{13}{\WJ{Dar cel care va răbda până la sfârșit va fi mântuit.
}}

\VS{14}{\WJ{Această Veste Bună a Împărăției va fi propovăduită în toată lumea, ca mărturie pentru toate națiunile, și atunci va veni sfârșitul.
}}


\par }{\PP \VS{15}“{\WJ{Când veți vedea urâciunea pustiirii, despre care s-a vorbit prin proorocul Daniel, stând în sfântul locaș (să înțeleagă cititorul),
}}

\VS{16}{\WJ{atunci cei ce sunt în Iudeea să fugă în munți.
}}

\VS{17}{\WJ{Cel care este pe acoperișul casei să nu se coboare să scoată lucrurile care sunt în casa lui.
}}

\VS{18}{\WJ{Cel care este pe câmp să nu se întoarcă să își ia hainele.
}}

\VS{19}{\WJ{Dar vai de cele însărcinate și de mamele care alăptează în acele zile!
}}

\VS{20}{\WJ{Rugați-vă ca fuga voastră să nu fie în timpul iernii și nici într-un Sabat,
}}

\VS{21}{\WJ{pentru că atunci va fi o mare suferință, așa cum nu a fost de la începutul lumii până acum, nu, și nici nu va fi vreodată.
}}

\VS{22}{\WJ{Dacă acele zile nu ar fi fost scurtate, nicio făptură nu ar fi fost salvată. Dar, de dragul celor aleși, acele zile vor fi scurtate.
}}


\par }{\PP \VS{23}{\WJ{Dacă vă va spune cineva: “Iată, aici este Hristosul!” sau: “Iată!”, nu-l credeți.
}}

\VS{24}Căci se
{\WJ{vor ridica falși hristoși și falși profeți, care vor face semne și minuni mari, ca să ducă în rătăcire, dacă este posibil, chiar și pe cei aleși.
}}


\par }{\PP \VS{25}“{\WJ{Iată, v-am spus mai dinainte.
}}


\par }{\PP \VS{26}{\WJ{De aceea, dacă vă vor spune: “Iată-l în pustie”, nu ieșiți; sau “Iată-l în lăcașuri”, nu credeți.
}}

\VS{27}{\WJ{Căci, după cum fulgerul fulgeră de la răsărit și se vede până la apus, așa va fi și venirea Fiului Omului.
}}

\VS{28}{\WJ{Căci oriunde este cadavrul, acolo se adună vulturii}}.


\par }{\PP \VS{29}“{\WJ{Dar îndată după suferința acelor zile, soarele se va întuneca, luna nu va mai lumina, stelele vor cădea din cer și puterile cerurilor se vor clătina;
}}

\VS{30}și
{\WJ{atunci va apărea pe cer semnul Fiului Omului. Atunci toate semințiile pământului vor plânge și vor vedea pe Fiul Omului venind pe norii cerului cu putere și cu mare slavă.
}}

\VS{31}{\WJ{El va trimite pe îngerii Săi cu un sunet mare de trâmbiță și vor aduna pe cei aleși ai Săi din cele patru vânturi, de la un capăt la altul al cerului.
}}


\par }{\PP \VS{32}{\WJ{“Învățați acum de la smochin pilda aceasta: Când ramura lui s-a înmuiat și produce frunze, știți că vara este aproape.
}}

\VS{33}Tot
{\WJ{așa și voi, când vedeți toate aceste lucruri, știți că El este aproape, chiar la uși.
}}

\VS{34}{\WJ{Adevărat vă spun că generația aceasta nu va trece până când nu se vor împlini toate aceste lucruri.
}}

\VS{35}{\WJ{Cerul și pământul vor trece, dar cuvintele Mele nu vor trece.
}}


\par }{\PP \VS{36}{\WJ{Dar nimeni nu știe ziua și ceasul acela, nici îngerii din ceruri, ci numai Tatăl Meu.
}}

\VS{37}{\WJ{Cum au fost zilele lui Noe, așa va fi și venirea Fiului Omului.
}}

\VS{38}{\WJ{Căci, după cum în zilele de dinaintea potopului mâncau și beau, se căsătoreau și se dădeau în căsătorie, până în ziua în care Noe a intrat în corabie,
}}

\VS{39}{\WJ{și nu au știut până când a venit potopul și i-a luat pe toți, așa va fi și venirea Fiului Omului.
}}

\VS{40}{\WJ{Atunci doi oameni vor fi pe câmp: unul va fi luat și unul va rămâne.
}}

\VS{41}{\WJ{Două femei vor fi măcinând la moară: una va fi luată și una va fi lăsată.
}}

\VS{42}{\WJ{Vegheați deci, căci nu știți în ce ceas vine Domnul vostru.
}}

\VS{43}{\WJ{Dar să știți că, dacă stăpânul casei ar fi știut în ce ceas al nopții va veni hoțul, ar fi vegheat și nu ar fi îngăduit să i se spargă casa.
}}

\VS{44}{\WJ{De aceea și voi fiți gata, căci în ceasul la care nu vă așteptați, Fiul Omului va veni.
}}


\par }{\PP \VS{45}{\WJ{Cine este robul credincios și înțelept, pe care l-a pus stăpânul său peste casa sa, ca să le dea hrana la vremea potrivită?
}}

\VS{46}{\WJ{Ferice de robul acela pe care îl găsește făcând așa când vine stăpânul său!
}}

\VS{47}Cu
{\WJ{siguranță vă spun că îl va pune peste tot ce are.
}}

\VS{48}{\WJ{Dar dacă robul acela rău va zice în inima lui: “Stăpânul meu întârzie să vină”,
}}

\VS{49}{\WJ{și va începe să-și bată tovarășii de slujbă, să mănânce și să bea cu bețivii,
}}

\VS{50}{\WJ{stăpânul acelui rob va veni într-o zi când nu se va aștepta și într-un ceas când nu va ști,
}}

\VS{51}și-l
{\WJ{va tăia în bucăți și-i va rândui partea lui cu ipocriții. Acolo va fi plânsul și scrâșnitul dinților.
}}



\par }\Chap{25}{\PP \VerseOne{1}{\WJ{“Atunci Împărăția Cerurilor va fi ca zece fecioare care și-au luat lămpile și au ieșit în întâmpinarea mirelui.
}}

\VS{2}{\WJ{Cinci dintre ele erau nebune, iar cinci erau înțelepte.
}}

\VS{3}Cele care
{\WJ{erau nebune, când și-au luat lămpile, nu au luat ulei cu ele,
}}

\VS{4}{\WJ{dar cele înțelepte au luat ulei în vasele lor împreună cu lămpile lor.
}}

\VS{5}{\WJ{Și, pe când întârzia mirele, toate au adormit și au dormit.
}}

\VS{6}{\WJ{Dar la miezul nopții s-a auzit un strigăt: “Iată! Vine mirele! Ieșiți în întâmpinarea lui!
}}

\VS{7}{\WJ{Atunci toate acele fecioare s-au sculat și și-au aranjat lămpile.
}}

\VS{8}Cele
{\WJ{nebune au zis celor înțelepte: “Dați-ne puțin din untdelemnul vostru, căci lămpile noastre se sting.
}}

\VS{9}{\WJ{Dar înțeleptele au răspuns: “Și dacă nu este destul pentru noi și pentru voi? Mergeți mai degrabă la cei care vând și cumpărați pentru voi înșivă'.
}}

\VS{10}În
{\WJ{timp ce ele se duceau să cumpere, a venit mirele, și cele care erau pregătite au intrat cu el la ospățul de nuntă, și ușa s-a închis.
}}

\VS{11}{\WJ{După aceea au venit și celelalte fecioare, zicând: “Doamne, Doamne, deschide-ne!”.
}}

\VS{12}{\WJ{Dar El a răspuns: “Adevărat vă spun că nu vă cunosc.
}}

\VS{13}{\WJ{Vegheați, deci, pentru că nu știți nici ziua, nici ceasul în care va veni Fiul Omului.
}}


\par }{\PP \VS{14}{\WJ{Căci este ca un om care se duce în altă țară, care și-a chemat slugile sale și le-a încredințat bunurile sale.
}}

\VS{15}{\WJ{Unuia i-a dat cinci talanți, altuia doi, altuia unul, și altuia unul, fiecăruia după putința lui. Apoi și-a continuat călătoria.
}}

\VS{16}{\WJ{Imediat, cel care a primit cei cinci talanți s-a dus și a făcut schimb cu ei și a făcut alți cinci talanți.
}}

\VS{17}La
{\WJ{fel, și cel care a primit cei doi a mai câștigat încă doi.
}}

\VS{18}{\WJ{Dar cel care a primit un talant s-a dus, a săpat în pământ și a ascuns banii stăpânului său.
}}


\par }{\PP \VS{19}{\WJ{“După multă vreme, a venit stăpânul acelor robi și a făcut socoteli cu ei.
}}

\VS{20}{\WJ{Cel care primise cei cinci talanți a venit și a adus alți cinci talanți, zicând: “Doamne, mi-ai dat cinci talanți. Iată, am câștigat alți cinci talanți în plus față de ei'.
}}


\par }{\PP \VS{21}{\WJ{Domnul său i-a zis: “Bine, slugă bună și credincioasă. Ai fost credincios peste puține lucruri, te voi pune peste multe lucruri. Intră în bucuria stăpânului tău'.
}}


\par }{\PP \VS{22}{\WJ{Și cel ce a primit cei doi talanți a venit și a zis: “Doamne, mi-ai dat doi talanți. Iată, am mai câștigat încă doi talanți în plus față de ei'.
}}


\par }{\PP \VS{23}{\WJ{Domnul său i-a zis: “Bine, slugă bună și credincioasă. Ai fost credincios în privința câtorva lucruri. Eu te voi pune peste multe lucruri. Intră în bucuria stăpânului tău'.
}}


\par }{\PP \VS{24}Cel care
{\WJ{primise un talant a venit și a zis: “Doamne, te-am cunoscut că ești un om aspru, că seceri de unde n-ai semănat și aduni de unde n-ai împrăștiat.
}}

\VS{25}M-am
{\WJ{temut, m-am dus și am ascuns talantul tău în pământ. Iată că ai ceea ce este al tău'.
}}


\par }{\PP \VS{26}{\WJ{Dar stăpânul său i-a răspuns: “Slugă rea și leneșă! Știai că secer unde n-am semănat și adun unde n-am împrăștiat.
}}

\VS{27}{\WJ{Trebuia deci să fi depus banii mei la bancheri și, la venirea mea, să-i primesc înapoi pe ai mei cu dobândă.
}}

\VS{28}Luați
{\WJ{deci talantul de la el și dați-l celui care are zece talanți.
}}

\VS{29}{\WJ{Căci fiecăruia care are i se va da și va avea din belșug, dar celui care nu are i se va lua chiar și ceea ce are.
}}

\VS{30}Aruncați-l pe
{\WJ{robul nefolositor în întunericul de afară, unde va fi plâns și scrâșnit din dinți.
}}


\par }{\PP \VS{31}{\WJ{“Dar când va veni Fiul Omului în slava Sa și toți sfinții îngeri cu El, atunci va ședea pe tronul slavei Sale.
}}

\VS{32}{\WJ{Înaintea Lui se vor aduna toate neamurile și El le va despărți unele de altele, așa cum un păstor desparte oile de capre.
}}

\VS{33}{\WJ{El va pune oile la dreapta Sa, iar caprele la stânga.
}}

\VS{34}{\WJ{Atunci Regele le va spune celor de la dreapta sa: 'Veniți, binecuvântați de Tatăl meu, moșteniți Împărăția pregătită pentru voi de la întemeierea lumii;
}}

\VS{35}{\WJ{căci mie mi-a fost foame și voi mi-ați dat de mâncare. Mi-a fost sete și mi-ați dat să beau. Am fost străin și m-ați primit în casă.
}}

\VS{36}am
{\WJ{fost gol și m-ați îmbrăcat. Am fost bolnav și m-ați vizitat. Am fost în închisoare și ați venit la mine.
}}


\par }{\PP \VS{37}{\WJ{Atunci cei drepți Îi vor răspunde și vor zice: “Doamne, când Te-am văzut noi flămând și Te-am hrănit, sau însetat și Ți-am dat de băut?
}}

\VS{38}{\WJ{Când Te-am văzut străin și Te-am primit la noi, sau gol și Te-am îmbrăcat?
}}

\VS{39}{\WJ{Când Te-am văzut bolnav sau în închisoare și am venit la Tine?”.
}}


\par }{\PP \VS{40}{\WJ{Împăratul le va răspunde: “Adevărat vă spun că, pentru că ați făcut aceasta unuia dintre acești frați ai Mei cei mai mici, Mie Mi-ați făcut-o.
}}

\VS{41}{\WJ{Apoi va spune și celor de la stânga: 'Duceți-vă de la Mine, blestemaților, în focul cel veșnic, pregătit diavolului și îngerilor lui;
}}

\VS{42}{\WJ{pentru că am fost flămând și nu Mi-ați dat să mănânc; Mi-a fost sete și nu Mi-ați dat să beau;
}}

\VS{43}am
{\WJ{fost străin și nu M-ați găzduit; gol și nu M-ați îmbrăcat; bolnav și în temniță și nu M-ați vizitat'.
}}


\par }{\PP \VS{44}{\WJ{Atunci vor răspunde și ei, zicând: Doamne, când Te-am văzut noi flămând, sau însetat, sau străin, sau gol, sau bolnav, sau în temniță, și nu Te-am ajutat?
}}


\par }{\PP \VS{45}{\WJ{Atunci El le va răspunde și le va zice: “Adevărat vă spun că, dacă n-ați făcut lucrul acesta unuia dintre cei mai mici dintre aceștia, nici Mie nu Mi l-ați făcut.
}}

\VS{46}{\WJ{Aceștiavor pleca la pedeapsa veșnică, dar cei drepți la viața veșnică.”
}}



\par }\Chap{26}{\PP \VerseOne{1}După ce a isprăvit toate aceste cuvinte, Isus a zis ucenicilor Săi:

\VS{2}“{\WJ{Știți că peste două zile vine Paștele și Fiul Omului va fi dat spre răstignire.”
}}


\par }{\PP \VS{3}Atunci preoții cei mai de seamă, cărturarii și bătrânii poporului s-au adunat în curtea marelui preot, care se numea Caiafa.

\VS{4}S-au sfătuit împreună ca să prindă pe Isus prin înșelăciune și să-L omoare.

\VS{5}Dar ei au spus: “Nu în timpul sărbătorii, ca să nu se producă o răscoală în popor”.


\par }{\PP \VS{6}Pe când era Isus în Betania, în casa lui Simon leprosul,

\VS{7}a venit la El o femeie care avea un vas de alabastru cu mir de mare preț, și a turnat mir pe capul Lui, pe când ședea la masă.

\VS{8}Ucenicii Lui, văzând aceasta, s-au indignat și au zis: “De ce această risipă?

\VS{9}pentru că acest unguent ar fi putut fi vândut cu mult și dat săracilor.”


\par }{\PP \VS{10}Dar Isus, știind aceasta, le-a zis: “{\WJ{De ce tulburați pe femeie? Ea a făcut o faptă bună pentru mine.
}}

\VS{11}{\WJ{Pentru că voi aveți întotdeauna săraci cu voi, dar pe mine nu mă aveți întotdeauna.
}}

\VS{12}{\WJ{Căci, turnând acest mir pe trupul meu, ea a făcut-o ca să mă pregătească pentru înmormântare.
}}

\VS{13}{\WJ{Adevărat vă spun că, oriunde se va propovădui această Bună Vestire în toată lumea, se va vorbi și despre ceea ce a făcut această femeie, ca o amintire a ei.”
}}


\par }{\PP \VS{14}Atunci unul din cei doisprezece, numit Iuda Iscarioteanul, s-a dus la preoții cei mai de seamă

\VS{15}și le-a zis: “Cât vreți să-mi dați dacă vi-l predau?” Și au cântărit pentru el treizeci de arginți.

\VS{16}Din acel moment a căutat prilej să îl trădeze.


\par }{\PP \VS{17}În prima zi a azimilor, ucenicii au venit la Isus și I-au zis: “Unde vrei să-Ți pregătim Paștile?”


\par }{\PP \VS{18}Și a zis: “{\WJ{Du-te în cetate la un oarecare și spune-i: “Învățătorul zice: “Vremea Mea este aproape. Voi sărbători Paștele în casa ta împreună cu discipolii mei."”””.
}}


\par }{\PP \VS{19}Ucenicii au făcut cum le poruncise Isus și au pregătit Paștele.


\par }{\PP \VS{20}Când s-a făcut seară, El stătea la masă cu cei doisprezece ucenici.

\VS{21}În timp ce mâncau, a zis: “Cu
{\WJ{siguranță vă spun că unul dintre voi Mă va trăda.”
}}


\par }{\PP \VS{22}Ei erau foarte mâhniți, și fiecare dintre ei a început să-L întrebe: “Nu sunt eu, Doamne, nu-i așa?”


\par }{\PP \VS{23}El a răspuns: “{\WJ{Cel ce și-a înmuiat mâna cu mine în vasul acesta mă va trăda.
}}

\VS{24}{\WJ{Fiul Omului merge așa cum este scris despre El, dar vai de omul prin care va fi trădat Fiul Omului! Mai bine ar fi fost pentru acel om dacă nu s-ar fi născut.”
}}


\par }{\PP \VS{25}Iuda, care L-a trădat, a răspuns: “Nu sunt eu, nu-i așa, Rabbi?”

\par }{\PP El i-a spus:
{\WJ{“Tu ai spus-o”.
}}


\par }{\PP \VS{26}Pe când mâncau ei, Isus a luat o pâine, a mulțumit pentru ea și a frânt-o. A dat-o ucenicilor și a zis: “{\WJ{Luați, mâncați; acesta este trupul Meu.”
}}

\VS{27}A luat paharul, a mulțumit și le-a dat, zicând: “{\WJ{Beți-l toți,
}}

\VS{28}{\WJ{căci acesta este sângele Meu, al noului legământ, care se varsă pentru mulți, spre iertarea păcatelor.
}}

\VS{29}{\WJ{Dar vă spun că nu voi mai bea din acest fruct al viței de vie de acum înainte, până în ziua în care îl voi bea din nou cu voi în Împărăția Tatălui meu.”
}}


\par }{\PP \VS{30}După ce au cântat un imn, au ieșit pe Muntele Măslinilor.


\par }{\PP \VS{31}Atunci Isus le-a zis: “În
{\WJ{noaptea aceasta, toți vă veți poticni din pricina Mea, căci este scris: “Voi lovi pe păstor și oile turmei vor fi risipite.
}}

\VS{32}{\WJ{Dar, după ce voi fi înviat, voi merge înaintea voastră în Galileea.”
}}


\par }{\PP \VS{33}Dar Petru i-a răspuns: “Chiar dacă toți se vor poticni din pricina ta, eu nu mă voi poticni niciodată.”


\par }{\PP \VS{34}Isus i-a zis: “{\WJ{Adevărat îți spun că în seara aceasta, până nu va cânta cocoșul, te vei lepăda de Mine de trei ori.”
}}


\par }{\PP \VS{35}Petru I-a zis: “Chiar dacă ar fi să mor cu Tine, nu Mă voi lepăda de Tine.” Și toți ucenicii au spus la fel.


\par }{\PP \VS{36}Atunci Isus a venit cu ei la un loc numit Ghetsimani și a zis ucenicilor Săi:
{\WJ{“Stați aici, până ce Mă voi duce acolo să Mă rog.”
}}

\VS{37}A luat cu El pe Petru și pe cei doi fii ai lui Zebedei, și a început să se întristeze și să fie foarte tulburat.

\VS{38}Și le-a zis: “{\WJ{Sufletul Meu este foarte îndurerat, până la moarte. Rămâneți aici și vegheați cu mine”.
}}


\par }{\PP \VS{39}El a înaintat puțin, a căzut cu fața la pământ și se ruga, zicând: “{\WJ{Tatăl Meu, dacă este cu putință, să treacă de la Mine paharul acesta; dar nu ceea ce vreau eu, ci ceea ce vrei Tu.”
}}


\par }{\PP \VS{40}A venit la ucenici și, găsindu-i dormind, a zis lui Petru: “{\WJ{Ce, n-ați putut să vegheați cu Mine un ceas?
}}

\VS{41}{\WJ{Vegheați și rugați-vă, ca să nu intrați în ispită. Într-adevăr, duhul este binevoitor, dar carnea este slabă.”
}}


\par }{\PP \VS{42}Și iarăși S-a dus a doua oară și S-a rugat, zicând: “{\WJ{Tatăl Meu, dacă paharul acesta nu poate să treacă de la Mine, dacă nu-l beau, facă-se voia Ta.”
}}


\par }{\PP \VS{43}A venit iarăși și i-a găsit dormind, căci aveau ochii îngreunați.

\VS{44}I-a lăsat din nou, s-a dus și s-a rugat a treia oară, spunând aceleași cuvinte.

\VS{45}Apoi a venit la ucenicii Săi și le-a zis: “{\WJ{Încă mai dormiți și vă odihniți? Iată că s-a apropiat ceasul și Fiul Omului este dat în mâinile păcătoșilor.
}}

\VS{46}{\WJ{Ridicați-vă, să mergem. Iată, cel care mă trădează este aproape”.
}}


\par }{\PP \VS{47}Pe când vorbea El încă, iată că a venit Iuda, unul din cei doisprezece, și împreună cu el o mare mulțime de preoți de seamă și de bătrâni ai poporului, cu săbii și cu ciomege.

\VS{48}Cel care l-a trădat le dăduse un semn, zicând: “Pe cine sărut eu, acela este. Prindeți-l!”.

\VS{49}Imediat s-a apropiat de Isus, a zis: “Salutări, Rabi!” și l-a sărutat.


\par }{\PP \VS{50}Isus i-a zis:
{\WJ{“Prietene, ce cauți aici?”
}}

\par }{\PP Atunci au venit, au pus mâinile pe Isus și L-au luat.

\VS{51}Iată că unul dintre cei care erau cu Isus a întins mâna și a scos sabia, a lovit pe slujitorul marelui preot și i-a tăiat urechea.


\par }{\PP \VS{52}Atunci Isus i-a zis: “{\WJ{Pune-ți sabia la locul ei, căci toți cei ce iau sabia vor muri de sabie.
}}

\VS{53}{\WJ{Sau crezi că n-aș putea să-L rog pe Tatăl Meu și chiar acum să-mi trimită mai mult de douăsprezece legiuni de îngeri?
}}

\VS{54}{\WJ{Cum s-ar împlini atunci Scripturile care trebuie să fie așa?”
}}


\par }{\PP \VS{55}În ceasul acela, Isus a zis mulțimii: “Ați
{\WJ{ieșit voi ca un tâlhar, cu săbii și cu ciomege, ca să Mă prindeți? Eu stăteam în fiecare zi în templu, învățând, și voi nu m-ați arestat.
}}

\VS{56}{\WJ{Dar toate acestea s-au întâmplat pentru ca să se împlinească Scripturile profeților.”
}}

\par }{\PP Atunci toți ucenicii l-au lăsat și au fugit.


\par }{\PP \VS{57}Cei ce luaseră pe Isus L-au dus la Caiafa, marele preot, unde erau adunați cărturarii și bătrânii.

\VS{58}Dar Petru L-a urmărit de la distanță până la curtea marelui preot, a intrat și a șezut cu ofițerii, ca să vadă sfârșitul.


\par }{\PP \VS{59}Iar preoții cei mai de seamă, bătrânii și tot sfatul căutau mărturii false împotriva lui Isus, ca să-L omoare,

\VS{60}dar n-au găsit niciuna. Chiar dacă s-au prezentat mulți martori mincinoși, nu au găsit niciunul. Dar, în cele din urmă, doi martori mincinoși au ieșit în față

\VS{61}și au spus: “Omul acesta a spus: “Pot să distrug templul lui Dumnezeu și să-l zidesc în trei zile””.


\par }{\PP \VS{62}Marele preot s-a ridicat în picioare și i-a zis: “Nu ai răspuns? Ce este ceea ce mărturisesc aceștia împotriva ta?”

\VS{63}Dar Isus a tăcut. Marele preot i-a răspuns: “Te conjur, pe Dumnezeul cel viu, să ne spui dacă Tu ești Hristosul, Fiul lui Dumnezeu.”


\par }{\PP \VS{64}Isus i-a zis: “{\WJ{Așa ai spus. Cu toate acestea, îți spun că, după aceea, îl vei vedea pe Fiul Omului șezând la dreapta Puterii și venind pe norii cerului.”
}}


\par }{\PP \VS{65}Atunci marele preot și-a rupt hainele și a zis: “A spus o blasfemie! De ce mai avem nevoie de alți martori? Iată, acum ați auzit blasfemia lui.

\VS{66}Ce părere aveți?”

\par }{\PP Ei au răspuns: “Este vrednic de moarte!”

\VS{67}Atunci l-au scuipat în față și l-au bătut cu pumnii, iar unii l-au pălmuit,

\VS{68}zicând: “Propovăduiește-ne, Hristoase! Cine te-a lovit?”.


\par }{\PP \VS{69}Petru ședea afară, în curte, și o slujnică a venit la el și i-a zis: “Și tu ai fost cu Iisus Galileanul!”


\par }{\PP \VS{70}Dar el a negat în fața tuturor, zicând: “Nu știu despre ce vorbiți.”


\par }{\PP \VS{71}După ce a ieșit afară pe prispă, l-a văzut altcineva și a zis celor ce erau acolo: “Și acesta a fost cu Isus din Nazaret.”


\par }{\PP \VS{72}El a negat din nou cu jurământ: “Nu-l cunosc pe omul acesta.”


\par }{\PP \VS{73}După puțină vreme, cei ce stăteau de față au venit și au zis lui Petru: “Cu siguranță și tu ești unul dintre ei, căci vorbirea ta te face cunoscut.”


\par }{\PP \VS{74}Atunci a început să înjure și să jure: “Nu-l cunosc pe omul acela!”

\par }{\PP Imediat a cântat cocoșul.

\VS{75}Petru și-a adus aminte de cuvântul pe care i-l spusese Isus:
{\WJ{“Înainte de a cânta cocoșul, te vei lepăda de Mine de trei ori”. Atunci a ieșit afară și a plâns cu amar.
}}



\par }\Chap{27}{\PP \VerseOne{1}Când a venit dimineața, toți preoții cei mai de seamă și bătrânii poporului s-au sfătuit împotriva lui Isus ca să-L omoare.

\VS{2}L-au legat, L-au dus și L-au predat guvernatorului Ponțiu Pilat.


\par }{\PP \VS{3}Atunci Iuda, care L-a trădat, văzând că Isus a fost osândit, a avut remușcări și a adus înapoi cei treizeci de arginți preoților cei mai de seamă și bătrânilor,

\VS{4}zicând: “Am păcătuit că am trădat sânge nevinovat.”

\par }{\PP Dar ei au spus: “Ce ne pasă nouă de asta? Voi vă ocupați de asta”.


\par }{\PP \VS{5}A aruncat piesele de argint în sanctuar și a plecat. Apoi s-a dus și s-a spânzurat.


\par }{\PP \VS{6}Preoții cei mai de seamă au luat arginții și au zis: “Nu este îngăduit să-i punem în vistierie, fiindcă este prețul sângelui.”

\VS{7}S-au sfătuit și au cumpărat cu ele câmpul olarului pentru a îngropa străinii.

\VS{8}De aceea, câmpul acela a fost numit până în ziua de azi “Câmpul de sânge”.

\VS{9}Atunci s-a împlinit ceea ce fusese spus prin profetul Ieremia, care zicea

\par }{\Q “Au luat cei treizeci de arginți,

\par }{\QB prețul celui căruia i s-a pus un preț,

\par }{\QB pe care unii dintre copiii lui Israel l-au taxat,


\par }{\Q \VS{10}și le-au dat pentru câmpul olarului,

\par }{\QB așa cum mi-a poruncit Domnul.”


\par }{\PP \VS{11}Isus a stat în picioare înaintea guvernatorului; și guvernatorul L-a întrebat: “Ești Tu Împăratul iudeilor?”

\par }{\PP Isus i-a zis:
{\WJ{“Așa zici tu”.
}}


\par }{\PP \VS{12}Când a fost acuzat de preoții cei mai de seamă și de bătrâni, n-a răspuns nimic.

\VS{13}Atunci Pilat i-a zis: “Nu auzi câte mărturisiri se fac împotriva ta?”


\par }{\PP \VS{14}El nu i-a răspuns, nici măcar un cuvânt, așa că guvernatorul s-a mirat foarte mult.


\par }{\PP \VS{15}La sărbătoare, guvernatorul obișnuia să elibereze mulțimii câte un prizonier pe care îl dorea.

\VS{16}Ei aveau atunci un prizonier de seamă, numit Baraba.

\VS{17}Când s-au adunat deci, Pilat le-a zis: “Pe cine vreți să vă eliberez? Pe Baraba sau pe Isus, care se numește Hristos?”

\VS{18}Căci știa că din cauza invidiei îl predaseră.


\par }{\PP \VS{19}Pe când ședea el pe scaunul de judecată, soția lui a trimis la el, zicând: “Nu te mai lega de acest neprihănit, căci astăzi am suferit multe în vis din pricina lui.”


\par }{\PP \VS{20}Iar preoții cei mai de seamă și bătrânii au convins mulțimile să ceară pe Baraba și să nimicească pe Isus.

\VS{21}Dar guvernatorul le-a răspuns: “Pe care dintre cei doi vreți să vi-l eliberez?”

\par }{\PP Ei au spus: “Baraba!”


\par }{\PP \VS{22}Pilat le-a zis: “Ce să fac lui Isus, care se numește Hristos?”

\par }{\PP Și toți i-au zis: “Să fie răstignit!”


\par }{\PP \VS{23}Dar guvernatorul a zis: “De ce? Ce rău a făcut?”

\par }{\PP Dar ei strigau foarte tare, zicând: “Să fie răstignit!”


\par }{\PP \VS{24}Pilat, văzând că nu se câștiga nimic, ci că se făcea tulburare, a luat apă și s-a spălat pe mâini în fața mulțimii, zicând: “Eu sunt nevinovat de sângele acestui neprihănit. Aveți grijă de el”.


\par }{\PP \VS{25}Și tot poporul a răspuns: “Sângele lui să fie peste noi și peste copiii noștri!”


\par }{\PP \VS{26}Apoi le-a dat drumul lui Baraba, iar pe Isus l-a biciuit și l-a dat spre răstignire.


\par }{\PP \VS{27}Atunci soldații guvernatorului au dus pe Isus în pretoriu și au adunat toată garnizoana împotriva Lui.

\VS{28}L-au dezbrăcat și I-au pus o haină stacojie.

\VS{29}I-au împletit o cunună de spini și i-au pus-o pe cap, iar în mâna dreaptă o trestie; au îngenuncheat înaintea lui și l-au batjocorit, zicând: “Slavă, Împăratul iudeilor!”

\VS{30}L-au scuipat, au luat trestia și l-au lovit în cap.

\VS{31}După ce l-au batjocorit, i-au luat haina, i-au pus hainele pe el și l-au dus ca să-l răstignească.


\par }{\PP \VS{32}Când au ieșit afară, au găsit un om din Cirene, numit Simon, și l-au silit să meargă cu ei, ca să-și ducă crucea.

\VS{33}Când au ajuns la un loc numit “Golgota”, adică “Locul craniului”,

\VS{34}i-au dat să bea vin acru amestecat cu fiere. După ce a gustat, nu a vrut să bea.

\VS{35}După ce l-au răstignit, i-au împărțit hainele, aruncând la sorți,

\VS{36}și, stând acolo, îl priveau.

\VS{37}I-au pus deasupra capului său acuzația scrisă împotriva lui: “ACESTA ESTE IISUS, REGELE IUDEILOR”.


\par }{\PP \VS{38}Și erau doi tâlhari răstigniți împreună cu El, unul la dreapta Lui și altul la stânga.


\par }{\PP \VS{39}Cei ce treceau pe acolo Îl ocărau, clătinându-și capetele

\VS{40}și zicând: “Tu, care distrugi templul și-l zidești în trei zile, mântuiește-te! Dacă ești Fiul lui Dumnezeu, coboară-te de pe cruce!”


\par }{\PP \VS{41}Preoții cei mai de seamă, de asemenea, batjocorind împreună cu cărturarii, cu fariseii și cu bătrânii, au zis:

\VS{42}“A mântuit pe alții, dar nu se poate mântui pe sine însuși. Dacă el este Regele lui Israel, să se dea jos de pe cruce acum și vom crede în el.

\VS{43}El se încrede în Dumnezeu. Să-l elibereze Dumnezeu acum, dacă îl vrea, pentru că a spus: “Eu sunt Fiul lui Dumnezeu”.”

\VS{44}Și tâlharii care au fost răstigniți împreună cu El au aruncat asupra Lui aceeași ocară.


\par }{\PP \VS{45}De la ceasul al șaselea a fost întuneric peste toată țara până la ceasul al nouălea.

\VS{46}Pe la ceasul al nouălea, Isus a strigat cu glas tare:
{\WJ{“Eli, Eli, lima sabachthani?”
}}Adică:
{\WJ{“Dumnezeul meu, Dumnezeul meu, de ce m-ai părăsit?”.
}}


\par }{\PP \VS{47}Unii dintre cei ce stăteau acolo, auzind, au zis: “Omul acesta cheamă pe Ilie.”


\par }{\PP \VS{48}Și îndată unul dintre ei a alergat, a luat un burete, l-a umplut cu oțet, l-a pus pe o trestie și i-a dat să bea.

\VS{49}Ceilalți au zis: “Lasă-l în pace. Să vedem dacă vine Ilie să-l salveze”.


\par }{\PP \VS{50}Isus a strigat din nou cu glas tare și și-a dat duhul.


\par }{\PP \VS{51}Și iată că perdeaua templului s-a rupt în două, de sus până jos. Pământul s-a cutremurat și stâncile s-au despicat.

\VS{52}Mormintele s-au deschis și multe trupuri ale sfinților care adormiseră au înviat;

\VS{53}și, ieșind din morminte după înviere, au intrat în cetatea sfântă și s-au arătat multora.


\par }{\PP \VS{54}Iar centurionul și cei ce erau cu el, care priveau pe Isus, văzând cutremurul și cele ce se întâmplau, s-au îngrozit și au zis: “Cu adevărat, acesta era Fiul lui Dumnezeu!”


\par }{\PP \VS{55}Și erau acolo multe femei care priveau de departe, și care Îl urmau pe Isus din Galileea și Îl slujeau.

\VS{56}Printre ele erau Maria Magdalena, Maria, mama lui Iacov și a lui Iose și mama fiilor lui Zebedei.


\par }{\PP \VS{57}Când s-a făcut seară, a venit un om bogat din Arimateea, numit Iosif, care era și el ucenic al lui Isus.

\VS{58}Acest om s-a dus la Pilat și a cerut trupul lui Isus. Atunci Pilat a poruncit ca trupul să fie predat.

\VS{59}Iosif a luat trupul, l-a înfășurat într-o pânză de in curată

\VS{60}și l-a așezat în mormântul său nou, pe care îl tăiase în stâncă. Apoi a rostogolit o piatră mare la ușa mormântului și a plecat.

\VS{61}Maria Magdalena era acolo, împreună cu cealaltă Maria, așezată în fața mormântului.


\par }{\PP \VS{62}A doua zi, adică a doua zi după ziua pregătirii, s-au adunat la Pilat preoții cei mai de seamă și fariseii, și au zis:

\VS{63}“Domnule, ne aducem aminte ce a spus înșelătorul acela, pe când era încă în viață: “După trei zile voi învia”.

\VS{64}Poruncește, așadar, să se asigure mormântul până a treia zi, ca nu cumva discipolii lui să vină noaptea și să-l fure și să spună poporului: “A înviat din morți”; și ultima înșelăciune va fi mai rea decât prima”.”


\par }{\PP \VS{65}Pilat le-a zis: “Aveți un paznic. Duceți-vă și asigurați-o cât mai bine”.

\VS{66}Ei s-au dus deci cu garda și au securizat mormântul, sigilând piatra.



\par }\Chap{28}{\PP \VerseOne{1}După Sabat, când a început să se lumineze de ziuă, în prima zi a săptămânii, Maria Magdalena și cealaltă Marie au venit să vadă mormântul.

\VS{2}Și iată că s-a făcut un cutremur mare, căci un înger al Domnului s-a coborât din cer și a venit și a rostogolit piatra de la ușă și a șezut pe ea.

\VS{3}Înfățișarea lui era ca un fulger, iar îmbrăcămintea lui era albă ca zăpada.

\VS{4}De frica lui, gărzile s-au cutremurat și s-au făcut ca niște oameni morți.

\VS{5}Îngerul le-a răspuns femeilor: “Nu vă temeți, căci știu că voi căutați pe Isus, care a fost răstignit.

\VS{6}El nu este aici, pentru că a înviat, așa cum a spus. Veniți să vedeți locul unde zăcea Domnul.

\VS{7}Mergeți repede și spuneți-le discipolilor Lui: “A înviat din morți și iată că merge înaintea voastră în Galileea; acolo Îl veți vedea. Iată, v-am spus”.


\par }{\PP \VS{8}Și au ieșit repede din mormânt, cu frică și cu mare bucurie, și au alergat să aducă vestea ucenicilor Lui.

\VS{9}Pe când se duceau să anunțe pe ucenicii Lui, iată că Isus le-a ieșit în întâmpinare și le-a zis:
{\WJ{“Bucurați-vă!”
}}

\par }{\PP Au venit, s-au prins de picioarele Lui și I s-au închinat.


\par }{\PP \VS{10}Atunci Isus le-a zis: “{\WJ{Nu vă temeți. Duceți-vă și spuneți fraților mei să meargă în Galileea și acolo Mă vor vedea.”
}}


\par }{\PP \VS{11}Pe când se duceau ei, iată că au venit în cetate niște păzitori și au spus preoților cei mai de seamă tot ce se întâmplase.

\VS{12}După ce s-au adunat cu bătrânii și s-au sfătuit, au dat soldaților o sumă mare de argint,

\VS{13}zicând: “Spuneți că discipolii lui au venit noaptea și l-au furat în timp ce noi dormeam.

\VS{14}Dacă acest lucru ajunge la urechile guvernatorului, îl vom convinge și vă vom scăpa de griji.”

\VS{15}Ei au luat banii și au făcut cum li s-a spus. Această zicală s-a răspândit printre evrei și continuă până astăzi.


\par }{\PP \VS{16}Dar cei unsprezece ucenici s-au dus în Galileea, pe muntele unde îi trimisese Isus.

\VS{17}Când L-au văzut, s-au închinat înaintea Lui; dar unii se îndoiau.

\VS{18}Isus a venit la ei și le-a vorbit, zicând: “{\WJ{Mie Mi s-a dat toată autoritatea în cer și pe pământ.
}}

\VS{19}{\WJ{Mergeți și faceți discipoli din toate neamurile, botezându-i în numele Tatălui și al Fiului și al Sfântului Duh,
}}

\VS{20}{\WJ{învățându-isă păzească toate lucrurile pe care vi le-am poruncit. Iată, Eu sunt cu voi în toate zilele, până la sfârșitul veacului.”}} Amin.


\par }